\chapter{Groups, second encounter}
\label{chap:groups2}

In this chapter we return to $\categ{Grp}$ and study several topics of
a less `general' nature than those considered in
\Cref{chap:groups}\@. Most of what we do here will apply exclusively
to \emph{finite} groups; this is an important example in its own
right, as it has spectacular applications (for example, in Galois
theory; cf.\
\Cref{sec:fields:short-march-through-applications-of-galois-theory}),
and it is a good subject from the expository point of view, since it
gives us the opportunity to see several general concepts at work in a
context that is complex enough to carry substance, but simple enough
(in this tiny selection of elementary topics) to be appreciated
easily.

\section{The conjugation action}
\label{sec:groups2:the-conjugation-action}


\subsection{Actions of groups on sets, reminder}
\label{subsec:groups2:actions-of-groups-on-sets-reminder}
Groups really shine when you let them act on something. This section
will make this point very effectively, since we will get surprisingly
precise results on finite groups by extremely simple-minded
applications of the elementary facts concerning group actions that we
established back in~\Cref{sec:groups:group-actions}.

Recall that we proved
(\Cref{prop:groups:transitive-action-is-g-mod-h}) that every
\emph{transitive} (left-)~action of a group $G$ on a set $S$ is, up to
a natural notion of isomorphism, `left-multiplication on the set of
left-cosets $G/H$'. Here, $H$ may be taken to be the stabilizer
$\operatorname{Stab}_G(a)$ of any element $a \in S$, that is
(\Cref{defn:groups:stabilizer}) the subgroup of $G$ fixing $a$. This
fact applies to the orbits of every left-action of $G$ on a set; in
particular, the number of elements in a finite orbit $O$ equals the
index of the stabilizer of any $a \in O$; in particular
(\Cref{coro:groups:orbit-divides-order}) the number of elements $|O|$
of an orbit must divide the order $|G|$ of $G$, if $G$ is finite.

These considerations may be packaged into a useful `counting' formula,
which we could call the \emph{class formula} for that action; this
name is usually reserved to the particular case of the action of $G$
onto itself by conjugation, which we will explore more carefully
below.

In order to state the formula, assume $G$ acts on a set $S$; for
$a \in S$, let $G_a$ denote the stabilizer
$\operatorname{Stab}_G(a)$. Also, let $Z$ be the set of \emph{fixed
  points} of the action:
\[
  Z = \{a \in S \mid (\forall g \in G) : ga = a\}.
\]
Note that $a \in Z \iff G_a = G$; we could say that $a \in Z$ if and
only if the orbit of $a$ is `trivial', in the sense that it consists
of $a$ alone.

\begin{prop}{}{prop:groups2:class-formula-action}
  Let $S$ be a finite set, and let $G$ be a group acting on $S$. With
  notation as above,
  \[
    |S| = |Z| + \sum_{a \in A} [G : G_a],
  \]
  where $A \subseteq S$ has exactly one element for each nontrivial
  orbit of the action.
\end{prop}
\begin{proof}
  The orbits form a partition of $S$, and $Z$ collects the trivial
  orbits; hence
  \[
    |S| = |Z| + \sum_{a \in A} |O_a|,
  \]
  where $O_a$ denotes the orbit of $a$. By
  \Cref{prop:groups:transitive-action-is-g-mod-h}, the order $|O_a|$
  equals the index of the stabilizer of $a$, yielding the statement.
\end{proof}

The main strength of \cref{prop:groups2:class-formula-action} rests in
the fact that, if $G$ is finite, each summand $[G : G_a]$
\emph{divides} the order of $G$ (and is $> 1$). This can be a strong
constraint, when some information is known about $|G|$. For example,
let's see what this says when $G$ is a $p$-group:

\begin{defn}{$p$-group}{defn:groups2:p-group}
  A \textbf{$p$-group} is a finite group whose order is a power of a
  prime integer $p$.
\end{defn}

\begin{coro}{}{coro:groups2:p-group-fixed-pts}
  Let $G$ be a $p$-group acting on a finite set $S$, and let $Z$ be
  the fixed point set of the action. Then
  \[
    |Z| \equiv |S| \pmod{p}.
  \]
\end{coro}
\begin{proof}
  Indeed, each summand $[G : G_a]$ in
  \cref{prop:groups2:class-formula-action} is a number larger than
  $1$, and a power of $p$; hence it is $0 \pmod{p}$.
\end{proof}

For instance, in certain situations this can be used to
establish\footnote{In this sense,
  \cref{prop:groups2:class-formula-action} is an instance of a class
  of results known as `fixed point theorems'. The reader will likely
  encounter a few such theorems in topology courses, where the role of
  the `size' of a set may be played by (for example) the Euler
  characteristic of a topological space.} that $Z \neq \emptyset$: see
\Cref{exer:groups2:p-group-action-has-fixed-points}. Such immediate
consequences of \cref{prop:groups2:class-formula-action} will assist
us below, in the proof of Sylow's theorems.

\subsection{Center, centralizer, conjugacy classes}
\label{subsec:groups2:center-centralizer-conjugacy-classes}
Recall (\Cref{exam:groups:action-by-left-multiplication-conjugation})
that every group $G$ acts on itself in at least two interesting ways:
by (left-) multiplication and by \emph{conjugation}. The latter action
is defined by the following $\rho : G \times G \to G$:
\[
  \rho(g, a) = gag^{-1}.
\]
As we know (\Cref{subsec:groups:actions-on-sets}), this datum is
equivalent to the datum of a certain group homomorphism:
\[
  \sigma : G \to S_G
\]
from $G$ to the permutation group on $G$.

This action highlights several interesting objects:

\begin{defn}{}{defn:groups2:center}
  The \emph{center} of $G$, denoted $Z(G)$, is the subgroup
  $\ker\sigma$ of $G$.
\end{defn}

Concretely, the center\footnote{Why `Z'? `Center' is \emph{Zentrum}
  auf Deutsch.} of $G$ is
\[
  Z(G) = \{g \in G \mid (\forall a \in G) : ga = ag\}.
\]
Indeed, $\sigma(g)$ is the identity in $S_G$ if and only if
$\sigma(g)$ acts as the identity on $G$; that is, if and only if
$gag^{-1} = a$ for all $a \in G$; that is, if and only if $g$ commutes
with all elements of $G$. In other words, the center is the set of
fixed points in $G$ under the conjugation action.

Note that the center of a group $G$ is automatically normal in $G$:
this is nearly immediate to check `by hand', but there is no need to
do so since it is a kernel by definition and kernels are normal.

A group $G$ is commutative if and only if $Z(G) = G$, that is, if and
only if the conjugation action is trivial on $G$.

In general, feel happy when you discover that the center of a group is
not trivial: this will often allow you to set up proofs by induction
on the number of elements of the group, by mod-ing out by the center
(this is, roughly, how we will prove the first Sylow theorem). Or note
the following useful fact, which comes in handy when trying to prove
that a group is commutative:

\begin{lem}{}{lem:groups2:quotient-center-cyclic}
  Let $G$ be a finite group, and assume $G/Z(G)$ is cyclic. Then $G$
  is commutative (and hence $G/Z(G)$ is in fact trivial).
\end{lem}
\begin{proof}
  (Cf.\ \Cref{exer:groups2:quot-ctr-iso-inn}.) As $G/Z(G)$ is cyclic,
  there exists an element $g \in G$ such that the class $gZ(G)$
  generates $G/Z(G)$. Then $\forall a \in G$
  \[
    aZ(G) = (gZ(G))^r
  \]
  for some $r \in \Z$; that is, there is an element $z \in Z(G)$ of
  the center such that $a = g^r z$.

  If now $a$, $b$ are in $G$, use this fact to write
  \[
    a = g^r z, \qquad b = g^s w
  \]
  for some $s \in \Z$ and $w \in Z(G)$; but then
  \[
    ab = (g^r z)(g^s w) = g^{r+s} zw = (g^s w)(g^r z) = ba,
  \]
  where I have used the fact that $z$ and $w$ commute with every
  element of $G$. As $a$ and $b$ were arbitrary, this proves that $G$
  is commutative.
\end{proof}

Next, the stabilizer of $a \in G$ under conjugation has a special
name:

\begin{defn}{}{defn:groups2:centralizer}
  The \emph{centralizer} (or \emph{normalizer}) $Z_G(a)$ of $a \in G$
  is its stabilizer under conjugation.
\end{defn}

Thus,
\[
  Z_G(a) = \{g \in G \mid gag^{-1} = a\} = \{g \in G \mid ga = ag\}
\]
consists of those elements in $G$ which commute with $a$. In
particular, $Z(G) \subseteq Z_G(a)$ for all $a \in G$; in fact,
$Z(G) = \bigcap_{a \in G} Z_G(a)$. Clearly
$a \in Z(G) \iff Z_G(a) = G$.

If there is no ambiguity concerning the group $G$ containing $a$, the
index $G$ may be dropped.

\begin{defn}{}{defn:groups2:conjugacy-class}
  The \emph{conjugacy class} of $a \in G$ is the orbit $[a]$ of $a$
  under the conjugation action. Two elements $a$, $b$ of $G$ are
  \emph{conjugate} if they belong to the same conjugacy class.
\end{defn}

The notation $[a]$ is not standard; $C(a)$ is used more frequently,
but I am not fond of it. Using $[a]$ reminds us that these are nothing
but the equivalence classes of elements of $G$ under a certain
interesting equivalence relation.

Note that $[a] = \{a\}$ if and only if $gag^{-1} = a$ for all
$g \in G$; that is, if and only if $ga = ag$ for all $g \in G$; that
is, if and only if $a \in Z(G)$.

\subsection{The Class Formula}
\label{subsec:groups2:the-class-formula}
The `official' Class Formula for a finite group $G$ is the particular
case of \cref{prop:groups2:class-formula-action} for the conjugation
action.

\begin{prop}{Class formula}{prop:groups2:class-formula}
  Let $G$ be a finite group. Then
  \[
    |G| = |Z(G)| + \sum_{a \in A} [G : Z(a)],
  \]
  where $A \subseteq G$ is a set containing one representative for
  each nontrivial conjugacy class in $G$.
\end{prop}
\begin{proof}
  The set of fixed points is $Z(G)$, and the stabilizer of $a$ is the
  centralizer $Z(a)$; apply \cref{prop:groups2:class-formula-action}.
\end{proof}
The class formula is surprisingly useful. In applying it, keep in mind
that every summand on the right (that is, both $|Z(G)|$ and each
$[G : Z(a)]$) is a divisor of $|G|$; this fact alone often suffices to
draw striking conclusions about $G$.

Possibly the most famous such application is to $p$-groups, via
\cref{coro:groups2:p-group-fixed-pts}:
\begin{coro}{}{coro:groups2:p-group-nontrivial-center}
  Let $G$ be a nontrivial $p$-group. Then $G$ has a nontrivial center.
\end{coro}
\begin{proof}
  Since $|Z(G)| \equiv |G| \bmod p$ and $|G| > 1$ is a power of $p$,
  necessarily $|Z(G)|$ is a multiple of $p$. As $Z(G) \neq \emptyset$
  (since $e_G \in Z(G)$), this implies $|Z(G)| \geq p$.
\end{proof}
For example, it follows immediately (from
\cref{coro:groups2:p-group-nontrivial-center} and
\cref{lem:groups2:quotient-center-cyclic}; cf.\
\Cref{exer:groups2:abel-or-triv-ctr-of-grp-ord-comp}) that if $p$ is
prime, then every group of order $p^2$ is commutative. In general, the
class formula poses a strong constraint on what can go on in a group.
\begin{exam}{}{exam:groups2:order-6-class-formula}
  Consider a group $G$ of order 6; what are the possibilities for its
  class formula?

  If $G$ is commutative, then the class formula will tell us very
  little:
  \[
    6 = 6.
  \]
  If $G$ is not commutative, then its center must be trivial (as a
  consequence of Lagrange's theorem and
  \cref{lem:groups2:quotient-center-cyclic}); so the class formula is
  $6 = 1 + \dots$, where $\dots$ collects the sizes of the nontrivial
  conjugacy classes. But each of these summands must be larger than 1,
  smaller than 6, and must divide 6; that is, there are no choices:
  \[
    6 = 1 + 2 + 3
  \]
  is the only possibility. The reader should check that this is indeed
  the class formula for $S_3$; in fact, $S_3$ is the only
  noncommutative group of order 6 up to isomorphism
  (\Cref{exer:groups2:noncomm-grp-ord-6-iso-s3}).
\end{exam}
Another useful observation is that normal subgroups must be unions of
conjugacy classes: because if $H$ is a normal subgroup, $a \in H$, and
$b = gag^{-1}$ is conjugate to $a$, then
\[
  b \in gHg^{-1} = H.
\]
To stick with the $|G| = 6$ example, note that every subgroup of a
group must contain the identity and its size must divide the order of
the group; it follows that a normal subgroup of a noncommutative group
of order 6 cannot have order 2, since 2 cannot be written as sums of
orders of conjugacy classes (including the class of the identity).

\subsection{Conjugation of subsets and subgroups}
\label{subsec:groups2:conjugation-of-subsets-and-subgroups}
We may also act by conjugation on the power set of $G$: if
$A \subseteq G$ is a subset and $g \in G$, the \emph{conjugate} of $A$
is the subset $gAg^{-1}$.  By cancellation, the conjugation map
$a \mapsto gag^{-1}$ is a bijection between $A$ and $gAg^{-1}$.

This leads to terminology analogous to the one introduced in
\Cref{subsec:groups2:center-centralizer-conjugacy-classes}.
\begin{defn}{}{defn:groups2:normalizer-centralizer-subset}
  The \emph{normalizer} $N_G(A)$ of $A$ is its stabilizer under
  conjugation.  The \emph{centralizer} of $A$ is the subgroup
  $Z_G(A) \subseteq N_G(A)$ fixing each element of $A$.
\end{defn}
Thus, $g \in N_G(A)$ if and only if\footnote{If $A$ is finite (but not
  in general), this condition is equivalent to
  $gAg^{-1} \subseteq A$.}  $gAg^{-1} = A$, and $g \in Z_G(A)$ if and
only if $\forall a \in A$, $gag^{-1} = a$.

For $A = \{a\}$ a singleton, we have
$N_G(\{a\}) = Z_G(\{a\}) = Z_G(a)$. In general,
$Z_G(A) \subsetneq N_G(A)$.

If $H$ is a subgroup of $G$, every conjugate $gHg^{-1}$ of $H$ is also
a subgroup of $G$; conjugate subgroups have the same order.
\begin{rem}{}{}
  The definition implies immediately that $H \subseteq N_G(H)$ and
  that $H$ is normal in $G$ if and only if $N_G(H) = G$. More
  generally, the normalizer $N_G(H)$ of $H$ in $G$ is (clearly) the
  largest subgroup of $G$ in which $H$ is normal.
\end{rem}
One could apply \cref{prop:groups2:class-formula-action} to the
conjugation action on subsets or subgroups; however, there are too
many subsets, and one has little control over the number of
subgroups. Other numerical considerations involving the number of
conjugates of a given subset or subgroup may be very useful.
\begin{lem}{}{lem:groups2:conjugate-subgroups-count}
  Let $H \subseteq G$ be a subgroup. Then (if finite) the number of
  subgroups conjugate to $H$ equals the index $[G : N_G(H)]$ of the
  normalizer of $H$ in $G$.
\end{lem}
\begin{proof}
  This is again an immediate consequence of
  \Cref{prop:groups:transitive-action-is-g-mod-h}.
\end{proof}
\begin{coro}{}{coro:groups2:conjugate-subgroups-divides-index}
  If $[G : H]$ is finite, then the number of subgroups conjugate to
  $H$ is finite and divides $[G : H]$.
\end{coro}
\begin{proof}
  \[ [G : H] = [G : N_G(H)] \cdot [N_G(H) : H]
  \]
  (cf.\ \Cref{subsec:groups:the-index-and-lagrange-s-theorem}).
\end{proof}
One of the celebrated Sylow theorems will strengthen this statement
substantially in the case in which $H$ is a maximal $p$-group
contained in a finite group $G$. For a statement concerning the size
of the normalizer of an arbitrary $p$-subgroup of a group, see
\Cref{lem:groups2:normalizer-index-mod-p}.  Another useful numerical
tool is the observation that if $H$ and $K$ are subgroups of a group
$G$ and $H \subseteq N_G(K)$---so that $gKg^{-1} = K$ for all
$g \in H$---then conjugation by $g \in H$ gives an automorphism of
$K$. Indeed, I have already observed that conjugation is a bijection,
and it is immediate to see that it is a homomorphism:
$\forall k_1, k_2 \in K$
\[
  (gk_1 g^{-1})(gk_2 g^{-1}) = gk_1(g^{-1}g)k_2 g^{-1} = g(k_1
  k_2)g^{-1}.
\]
Thus, conjugation gives a set-function
\[
  \gamma : H \to \Aut_{\categ{Grp}}(K).
\]
The reader will check that this is a group homomorphism and will
determine $\ker\gamma$
(\Cref{exer:groups2:conj-action-kernel-is-centralizer}).

This is especially useful if $H$ is finite and some information is
available concerning $\Aut_{\categ{Grp}}(K)$ (for an example, see
\Cref{exer:groups2:ctr-of-an-trivial}). A classic application is
presented in \Cref{exer:groups2:min-prime-ord-normal-subgrp-central}.

\subsection*{Exercises}

\begin{exer}{}{exer:groups2:p-group-action-has-fixed-points}
  \exerused{} Let $p$ be a prime integer, let $G$ be a $p$-group, and
  let $S$ be a set such that $|S| \not\equiv 0 \bmod p$. If $G$ acts
  on $S$, prove that the action must have fixed
  points. [\Cref{subsec:groups2:actions-of-groups-on-sets-reminder},
  \Cref{subsec:groups2:sylow-ii}]
\end{exer}
\begin{exer}{}{exer:groups2:ctr-dhdrl-grp-of-present}
  Find the center of $D_{2n}$. (The answer depends on the parity of
  $n$. You have actually done this already:
  \Cref{exer:groups:dihedral-center}. This time, use a presentation.)
\end{exer}
\begin{exer}{}{exer:groups2:ctr-symm-grp-trivial}
  Prove that the center of $S_n$ is trivial for $n \geq 3$. (Suppose
  that $\sigma \in S_n$ sends $a$ to $b \neq a$, and let
  $c \neq a, b$. Let $\tau$ be the permutation that acts solely by
  swapping $b$ and $c$. Then compare the action of $\sigma\tau$ and
  $\tau\sigma$ on $a$.)
\end{exer}
\begin{exer}{}{exer:groups2:subgrp-of-ctr-normal}
  \exerused{} Let $G$ be a group, and let $N$ be a subgroup of
  $Z(G)$. Prove that $N$ is normal in $G$. [\Cref{subsec:groups2:sylow-i}]
\end{exer}
\begin{exer}{}{exer:groups2:quot-ctr-iso-inn}
  \exerused{} Let $G$ be a group. Prove that $G/Z(G)$ is isomorphic to
  the group $\operatorname{Inn}(G)$ of inner automorphisms of
  $G$. (Cf.\ \Cref{exer:groups:inner-automorphisms}.) Then prove
  \cref{lem:groups2:quotient-center-cyclic} again by using the result
  of
  \Cref{exer:groups:inn-cyclic-iff-abelian}. [\Cref{subsec:groups2:center-centralizer-conjugacy-classes}]
\end{exer}
\begin{exer}{}{exer:groups2:abel-or-triv-ctr-of-grp-ord-comp}
  \exerused{} Let $p$, $q$ be prime integers, and let $G$ be a group
  of order $pq$. Prove that either $G$ is commutative or the center of
  $G$ is trivial.  Conclude (using
  \cref{coro:groups2:p-group-nontrivial-center}) that every group of
  order $p^2$, for a prime $p$, is commutative. [\Cref{subsec:groups2:the-class-formula}]
\end{exer}
\begin{exer}{}{exer:groups2:p3-abel-Q}
  Prove or disprove that if $p$ is prime, then every group of order
  $p^3$ is commutative.
\end{exer}
\begin{exer}{}{exer:groups2:p-grp-norm-subgrp-ord-pk}
  \exerused{} Let $p$ be a prime number, and let $G$ be a $p$-group:
  $|G| = p^r$.  Prove that $G$ contains a normal subgroup of order
  $p^k$ for every nonnegative $k \leq r$. [\Cref{subsec:groups2:sylow-i}]
\end{exer}
\begin{exer}{}{exer:groups2:nontriv-norm-subgrp-meets-ctr}
  $\lnot$ Let $p$ be a prime number, $G$ a $p$-group, and $H$ a
  nontrivial normal subgroup of $G$. Prove that
  $H \cap Z(G) \neq \{e\}$. (Hint: Use the class formula.)
  [\ref{exer:groups2:nilp-grp-normal-subgrp-meets-ctr}]
\end{exer}
\begin{exer}{}{exer:groups2:odd-ord-conj-inv-implies-triv}
  Prove that if $G$ is a group of odd order and $g \in G$ is conjugate
  to $g^{-1}$, then $g = e_G$.
\end{exer}
\begin{exer}{}{exer:groups2:commuting-class-reps-implies-abel}
  Let $G$ be a finite group, and suppose there exist representatives
  $g_1, \dots, g_r$ of the $r$ distinct conjugacy classes in $G$, such
  that $\forall i, j$, $g_i g_j = g_j g_i$. Prove that $G$ is
  commutative. (Hint: What can you say about the sizes of the
  conjugacy classes?)
\end{exer}
\begin{exer}{}{exer:groups2:class-formula-d8-q8-verify}
  Verify that the class formula for both $D_8$ and $Q_8$ (cf.\
  \Cref{exer:rings:quaternions}) is $8 = 2 + 2 + 2 + 2$. (Also note
  that $D_8 \not\cong Q_8$.)
\end{exer}
\begin{exer}{}{exer:groups2:noncomm-grp-ord-6-iso-s3}
  \exerused{} Let $G$ be a noncommutative group of order 6. As
  observed in \Cref{exam:groups2:order-6-class-formula}, $G$ must have
  trivial center and exactly two conjugacy classes, of order 2 and 3.
  \begin{itemize}
  \item Prove that if every element of a group has order $\leq 2$,
    then the group is commutative. Conclude that $G$ has an element
    $y$ of order 3.
  \item Prove that $\genalg{y}$ is normal in $G$.
  \item Prove that $[y]$ is the conjugacy class of order 2 and
    $[y] = \{y, y^2\}$.
  \item Prove that there is an $x \in G$ such that $yx = xy^2$.
  \item Prove that $x$ has order 2.
  \item Prove that $x$ and $y$ generate $G$.
  \item Prove that $G \cong S_3$.
  \end{itemize} [\Cref{subsec:groups2:the-class-formula}, \Cref{subsec:groups2:applications}]
\end{exer}
\begin{exer}{}{exer:groups2:conj-of-subset-bound-by-idx-ctr}
  Let $G$ be a group, and assume $[G : Z(G)] = n$ is finite. Let
  $A \subseteq G$ be any subset. Prove that the number of conjugates
  of $A$ is at most $n$.
\end{exer}
\begin{exer}{}{exer:groups2:class-formula-60-only-triv-normal}
  Suppose that the class formula for a group $G$ is
  $60 = 1 + 15 + 20 + 12 + 12$. Prove that the only normal subgroups
  of $G$ are $\{e\}$ and $G$.
\end{exer}
\begin{exer}{}{exer:groups2:idx-2-subgrp-conj-class-half}
  \exerused{} Let $G$ be a finite group, and let $H \subseteq G$ be a
  subgroup of index 2. For $a \in H$, denote by $[a]_H$, resp.,
  $[a]_G$, the conjugacy class of $a$ in $H$, resp., $G$. Prove that
  either $[a]_H = [a]_G$ or $[a]_H$ is half the size of $[a]_G$,
  according to whether the centralizer $Z_G(a)$ is not or is contained
  in $H$. (Hint: Note that $H$ is normal in $G$, by
  \Cref{exer:groups:index-2-subgroup-normal}; apply
  \Cref{prop:groups:second-isomorphism-theorem}.)
  [\Cref{subsec:groups2:conjugacy-an-simplicity}]
\end{exer}
\begin{exer}{}{exer:groups2:grp-not-union-of-conjugates}
  $\lnot$ Let $H$ be a proper subgroup of a finite group $G$. Prove
  that $G$ is not the union of the conjugates of $H$. (Hint: You know
  the number of conjugates of $H$; keep in mind that any two subgroups
  overlap, at least at the identity.)
  [\ref{exer:groups2:transitive-action-has-free-elt},
  \ref{exer:groups2:gl2c-union-of-conjugates-of-borel}]
\end{exer}
\begin{exer}{}{exer:groups2:transitive-action-has-free-elt}
  Let $S$ be a set endowed with a transitive action of a finite group
  $G$, and assume $|S| \geq 2$. Prove that there exists a $g \in G$
  without fixed points in $S$, that is, such that $gs \neq s$ for all
  $s \in S$. (Hint: By
  \Cref{prop:groups:transitive-action-is-g-mod-h}, you may assume
  $S = G/H$, with $H$ proper in $G$. Use
  \Cref{exer:groups2:grp-not-union-of-conjugates}.)
\end{exer}
\begin{exer}{}{exer:groups2:conj-class-disjoint-from-subgrp}
  Let $H$ be a proper subgroup of a finite group $G$. Prove that there
  exists a $g \in G$ whose conjugacy class is disjoint from $H$.
\end{exer}
\begin{exer}{}{exer:groups2:gl2c-union-of-conjugates-of-borel}
  Let $G = \mathrm{GL}_2(\C)$, and let $H$ be the subgroup consisting
  of upper triangular matrices
  (\Cref{exer:groups:upper-triangular-subgroup}). Prove that $G$ is
  the union of the conjugates of $H$. Thus, the finiteness hypothesis
  in \Cref{exer:groups2:grp-not-union-of-conjugates} is
  necessary. (Hint: Equivalently, prove that every $2 \times 2$ matrix
  is conjugate to a matrix in $H$. You will use the fact that $\C$ is
  algebraically closed; see
  \Cref{exam:rings:speck-prime-ideals-maximal}.)
\end{exer}
\begin{exer}{}{exer:groups2:conj-action-kernel-is-centralizer}
  \exerused{} Let $H$, $K$ be subgroups of a group $G$, with
  $H \subseteq N_G(K)$.  Verify that the function
  $\gamma : H \to \Aut_{\categ{Grp}}(K)$ defined by conjugation is a
  homomorphism of groups and that $\ker\gamma = H \cap Z_G(K)$, where
  $Z_G(K)$ is the centralizer of $K$.
  [\Cref{subsec:groups2:conjugation-of-subsets-and-subgroups},
  \ref{exer:groups2:min-prime-ord-normal-subgrp-central}]
\end{exer}
\begin{exer}{}{exer:groups2:min-prime-ord-normal-subgrp-central}
  \exerused{} Let $G$ be a finite group, and let $H$ be a cyclic
  subgroup of $G$ of order $p$. Assume that $p$ is the smallest prime
  dividing the order of $G$ and that $H$ is normal in $G$. Prove that
  $H$ is contained in the center of $G$. (Hint: By
  \Cref{exer:groups2:conj-action-kernel-is-centralizer} there is a
  homomorphism $\gamma : G \to \Aut_{\categ{Grp}}(H)$; by
  \Cref{exer:groups:aut-cn-euler-phi}, $\Aut_{\categ{Grp}}(H)$ has
  order $p - 1$. What can you say about $\gamma$?)
  [\Cref{subsec:groups2:conjugation-of-subsets-and-subgroups}]
\end{exer}

\section{The Sylow theorems}
\label{sec:groups2:the-sylow-theorems}


\subsection{Cauchy's theorem}
\label{subsec:groups2:cauchy-s-theorem}
The `Sylow theorems' consist of three statements concerning
$p$-subgroups (cf.\ \cref{defn:groups2:p-group}) of a given finite
group $G$. The form I will give for the first of these statements will
tell us that $G$ contains $p$-groups of all sizes allowed by
Lagrange's theorem: if $p$ is a prime and $p^k$ divides $|G|$, then
$G$ contains a subgroup of order $p^k$. The proof of this statement is
an easy induction, provided the statement for $k = 1$ is known: that
is, provided that one has established
\begin{thm}{Cauchy's theorem}{thm:groups2:cauchy}
  Let $G$ be a finite group, and let $p$ be a prime divisor of
  $|G|$. Then $G$ contains an element of order $p$.
\end{thm}
As it happens, only the \emph{abelian} version of this statement is
needed for the proof of the first Sylow theorem; then the full
statement of Cauchy's theorem follows from the first Sylow theorem
itself. Since the (diligent) reader has already proved Cauchy's
theorem for abelian groups (in \Cref{exer:groups:cauchy-abelian}), we
could directly move on to Sylow theorems.  However, there is a quick
proof\footnote{This argument is apparently due to James McKay.} of the
full statement of Cauchy's theorem which does not rely on Sylow and is
a good illustration of the power of the general `class formula for
arbitrary actions' (\cref{prop:groups2:class-formula-action}). I will
present this proof, while also encouraging the reader to go back and
(re)do \Cref{exer:groups:cauchy-abelian} now.
\begin{proof}[Proof of \cref{thm:groups2:cauchy}]
  Consider the set $S$ of ordered $p$-tuples of elements of $G$:
  \[
    (a_1, \dots, a_p)
  \]
  such that $a_1 \cdots a_p = e$. I claim that $|S| = |G|^{p-1}$:
  indeed, once $a_1, \dots, a_{p-1}$ are chosen (arbitrarily), then
  $a_p$ is determined as it is the inverse of $a_1 \cdots
  a_{p-1}$. Therefore, $p$ divides the order of $S$ as it divides the
  order of $G$.  Also note that if $a_1 \cdots a_p = e$, then
  \[
    a_2 \cdots a_p a_1 = e
  \]
  (even if $G$ is not commutative): because if $a_1$ is a left-inverse
  to $a_2 \cdots a_p$, then it is also a right-inverse to it.
  Therefore, we may act with the group $\Z/p\Z$ on $S$: given $[m]$ in
  $\Z/p\Z$, with $0 \leq m < p$, act by $[m]$ on $(a_1, \dots, a_p)$
  by sending it to
  \[
    (a_{m+1}, \dots, a_p, a_1, \dots, a_m)\,;
  \]
  as we just observed, this is still an element of $S$.  Now
  \cref{coro:groups2:p-group-fixed-pts} implies
  \[
    |Z| \equiv |S| \equiv 0 \bmod p,
  \]
  where $Z$ is the set of fixed points of this action. Fixed points
  are $p$-tuples of the form
  \[
    \tag{$*$} (a, \dots, a);
  \]
  and note that $Z \neq \emptyset$, since $\{e, \dots, e\} \in
  Z$. Since $p \geq 2$ and $p$ divides $|Z|$, we conclude that
  $|Z| > 1$; therefore there exists some element in $Z$ of the
  form~$(*)$, with $a \neq e$.

  This says that there exists an $a \in G$, $a \neq e$, such that
  $a^p = e$, proving the statement.
\end{proof}
We should remark that the proof given here proves a more precise
result than the raw statement of \cref{thm:groups2:cauchy}: every
element of order $p$ in $G$ generates a cyclic subgroup of $G$ of
order $p$, and we are able to say something about the number of such
subgroups.
\begin{fact}{}{fact:groups2:cauchy-cyclic-count}
  Let $G$ be a finite group, let $p$ be a prime divisor of $|G|$, and
  let $N$ be the number of cyclic subgroups of $G$ of order $p$. Then
  $N \equiv 1 \bmod p$.
\end{fact}

The proof of this fact is left to the reader (as an incentive to
really understand the proof of \cref{thm:groups2:cauchy}).
\cref{fact:groups2:cauchy-cyclic-count}, coupled with the simple
observation that if there is only 1 cyclic subgroup $H$ of order $p$,
then that subgroup must be normal
(\Cref{exer:groups2:characteristic-subgrp-basics}), suffices for
interesting applications.

\begin{defn}{}{defn:groups2:simple-group}
  A group $G$ is \emph{simple} if it is nontrivial and its only normal
  subgroups are $\{e\}$ and $G$ itself.
\end{defn}
Simple groups occupy a special place in the theory of groups: one can
`break up' any finite group into basic constituents which are simple
groups; we will see how this is done in
\Cref{subsec:groups2:the-jordan-h-older-theorem}. Thus, it is
important to be able to tell whether a group is simple or
not.\footnote{In fact, a complete list of all finite simple groups is
  known: this is the classification result mentioned at the end of
  \Cref{subsec:groups:example-subgroup-generated-by-a-subset},
  arguably one of the deepest and hardest results in mathematics.}
\begin{exam}{}{exam:groups2:mp-not-simple}
  Let $p$ be a positive prime integer. If $|G| = mp$, with
  $1 < m < p$, then $G$ is not simple.

  Indeed, consider the subgroups of $G$ with $p$ elements. By
  \cref{fact:groups2:cauchy-cyclic-count}, the number of such
  subgroups is $\equiv 1 \bmod p$. Thus, if there is more than one
  such subgroup, then there must be at least $p + 1$. Any two distinct
  subgroups of prime order can only meet at the identity (why?);
  therefore this would account for at least
  \[
    1 + (p+1)(p-1) = p^2
  \]
  elements in $G$. Since $|G| = mp < p^2$, this is
  impossible. Therefore there is only one cyclic subgroup of order $p$
  in $G$, which must be normal as mentioned above, proving that $G$ is
  not simple.
\end{exam}

\subsection{Sylow I}
\label{subsec:groups2:sylow-i}
Let $p$ be a prime integer. A \emph{$p$-Sylow subgroup} of a finite
group $G$ is a subgroup of order $p^r$, where $|G| = p^r m$ and
$\gcd(p, m) = 1$. That is, $P \subseteq G$ is a $p$-Sylow subgroup if
it is a $p$-group and $p$ does not divide $[G : P]$.

If $p$ does not divide the order of $G$, then $G$ contains a $p$-Sylow
subgroup: namely, $\{e\}$. This is not very interesting; what is
interesting is that $G$ contains a $p$-Sylow subgroup even when $p$
\emph{does} divide the order of $G$:
\begin{thm}{First Sylow theorem}{thm:groups2:sylow-i}
  Every finite group contains a $p$-Sylow subgroup, for all primes
  $p$.
\end{thm}

The first Sylow theorem follows from the seemingly stronger statement:
\begin{prop}{}{prop:groups2:sylow-i-strong}
  If $p^k$ divides the order of $G$, then $G$ has a subgroup of order
  $p^k$.
\end{prop}

The statements are actually easily seen to be equivalent, by
\cref{exer:groups2:p-grp-norm-subgrp-ord-pk}; in any case, the
standard argument proving \cref{thm:groups2:sylow-i} proves
\cref{prop:groups2:sylow-i-strong}, and I see no reason to hide this
fact. Here is the argument:
\begin{proof}[Proof of \cref{prop:groups2:sylow-i-strong}]
  If $k = 0$, there is nothing to prove, so we may assume $k \geq 1$
  and in particular that $|G|$ is a multiple of $p$.

  Argue by induction on $|G|$: if $|G| = p$, again there is nothing to
  prove; if $|G| > p$ and $G$ contains a proper subgroup $H$ such that
  $[G : H]$ is relatively prime to $p$, then $p^k$ divides the order
  of $H$, and hence $H$ contains a subgroup of order $p^k$ by the
  induction hypothesis, and thus so does $G$.  Therefore, we may
  assume that all proper subgroups of $G$ have index divisible by
  $p$. By the class formula (\cref{prop:groups2:class-formula}), $p$
  divides the order of the center $Z(G)$.

  By Cauchy's theorem\footnote{Note that, as mentioned in
    \Cref{subsec:groups2:cauchy-s-theorem}, we only need the
    abelian case of this theorem.}, $\exists\, a \in Z(G)$ such that
  $a$ has order $p$. The cyclic subgroup $N = \genalg{a}$ is contained
  in $Z(G)$, and hence it is normal in $G$
  (\cref{exer:groups2:subgrp-of-ctr-normal}). Therefore we can
  consider the quotient $G/N$.  Since $|G/N| = |G|/p$ and $p^k$
  divides $|G|$ by hypothesis, we have that $p^{k-1}$ divides the
  order of $G/N$. By the induction hypothesis, we may conclude that
  $G/N$ contains a subgroup of order $p^{k-1}$. By the structure of
  the subgroups of a quotient
  (\Cref{subsec:groups:subgroups-of-quotients}, especially
  \Cref{prop:groups:subgroups-of-quotients-bijection}), this subgroup
  must be of the form $P/N$, for $P$ a subgroup of $G$.

  But then $|P| = |P/N| \cdot |N| = p^{k-1} \cdot p = p^k$, as needed.
\end{proof}
There are slicker ways to prove \cref{thm:groups2:sylow-i}. We will
see a pretty (and insightful) alternative in
\Cref{subsec:groups2:sylow-ii}; but the proof given above is easy
to remember and is a good template for similar arguments.

\begin{rem}{}{}
  The diligent reader worked out in
  \Cref{exer:groups:abelian-subgroup-of-order-d} a \emph{stronger}
  statement than \cref{prop:groups2:sylow-i-strong}, for abelian
  groups. The arguments are similar; the advantage in the abelian case
  is that any cyclic subgroup produced by Cauchy's theorem is
  automatically normal, while ensuring normality requires a few twists
  and turns in the general case (and, as a result, yields a weaker
  statement).
\end{rem}

\subsection{Sylow II}
\label{subsec:groups2:sylow-ii}
\cref{thm:groups2:sylow-i} tells us that some maximal $p$-group in $G$
attains the largest size allowed by Lagrange's theorem, that is, the
maximal power of the prime $p$ dividing $|G|$.

One can be more precise: the second Sylow theorem tells us that every
maximal $p$-group in $G$ is in fact a $p$-Sylow subgroup. It is as
large as is allowed by Lagrange's theorem.

The situation is in fact even better: all $p$-Sylow subgroups are
conjugates of each other.\footnote{Of course if $P$ is a $p$-Sylow
  subgroup of $G$, then so are all conjugates $gPg^{-1}$ of $P$.}
Moreover, even better than this, every $p$-group inside $G$ must be
contained in a conjugate of any fixed $p$-Sylow subgroup.

The proof of this very precise result is very easy!
\begin{thm}{Second Sylow theorem}{thm:groups2:sylow-ii}
  Let $G$ be a finite group, let $P$ be a $p$-Sylow subgroup, and let
  $H \subseteq G$ be a $p$-group. Then $H$ is contained in a conjugate
  of $P$: there exists $g \in G$ such that $H \subseteq gPg^{-1}$.
\end{thm}
\begin{proof}
  Act with $H$ on the set of left-cosets of $P$, by
  left-multiplication. Since there are $[G : P]$ cosets and $p$ does
  not divide $[G : P]$, we know this action must have fixed points
  (\cref{exer:groups2:p-group-action-has-fixed-points}): let $gP$ be
  one of them. This means that $\forall h \in H$:
  \[
    hgP = gP;
  \]
  that is, $g^{-1}hgP = P$ for all $h$ in $H$; that is,
  $g^{-1}Hg \subseteq P$; that is, $H \subseteq gPg^{-1}$, as needed.
\end{proof}
We can obtain an even more complete picture of the situation. Suppose
we have constructed a chain
\[
  H_0 = \{e\} \subseteq H_1 \subseteq \dots \subseteq H_k
\]
of $p$-subgroups of a group $G$, where $|H_i| = p^i$. By
\cref{thm:groups2:sylow-ii} we know that $H_k$ is contained in some
$p$-Sylow subgroup, of order $p^r$ = the maximum power of $p$ dividing
the order of $G$. But I claim that the chain can in fact be continued
one step at a time all the way up to the Sylow subgroup:
\[
  H_0 = \{e\} \subseteq H_1 \subseteq \dots \subseteq H_k \subseteq
  H_{k+1} \subseteq \dots \subseteq H_r;
\]
and, further, $H_k$ may be assumed to be normal in $H_{k+1}$. The
following lemma will simplify the proof of this fact considerably and
will also help us prove the third Sylow theorem.
\begin{lem}{}{lem:groups2:normalizer-index-mod-p}
  Let $H$ be a $p$-group contained in a finite group $G$. Then
  \[ [N_G(H) : H] \equiv [G : H] \bmod p.
  \]
\end{lem}
\begin{proof}
  If $H$ is trivial, then $N_G(H) = G$ and the two numbers are equal.

  Assume then that $H$ is nontrivial, and act with $H$ on the set of
  left-cosets of $H$ in $G$, by left-multiplication. The fixed points
  of this action are the cosets $gH$ such that $\forall h \in H$
  \[
    hgH = gH,
  \]
  that is, such that $g^{-1}hg \in H$ for all $h \in H$; in other
  words, $H \subseteq gHg^{-1}$, and hence (by order considerations)
  $gHg^{-1} = H$.  This means precisely that $g \in
  N_G(H)$. Therefore, the set of fixed points of the action consists
  of the set of cosets of $H$ in $N_G(H)$.

  The statement then follows immediately from
  \cref{coro:groups2:p-group-fixed-pts}.
\end{proof}
As a consequence, if $H_k$ is not a $p$-Sylow subgroup `already', in
the sense that $p$ `still' divides $[G : H_k]$, then $p$ must also
divide $[N_G(H_k) : H_k]$. Another application of Cauchy's theorem
tells us how to obtain the next subgroup $H_{k+1}$ in the chain. More
precisely, we have the following result.
\begin{prop}{}{prop:groups2:p-subgroup-chain-step}
  Let $H$ be a $p$-subgroup of a finite group $G$, and assume that $H$
  is not a $p$-Sylow subgroup. Then there exists a $p$-subgroup $H'$
  of $G$ containing $H$, such that $[H' : H] = p$ and $H$ is normal in
  $H'$.
\end{prop}
\begin{proof}
  Since $H$ is not a $p$-Sylow subgroup of $G$, $p$ divides
  $[N_G(H) : H]$, by \cref{lem:groups2:normalizer-index-mod-p}. Since
  $H$ is normal in $N_G(H)$, we may consider the quotient group
  $N_G(H)/H$, and $p$ divides the order of this group. By
  \cref{thm:groups2:cauchy}, $N_G(H)/H$ has an element of order $p$;
  this generates a subgroup of order $p$ of $N_G(H)/H$, which must be
  (cf.\ \Cref{subsec:groups:subgroups-of-quotients}) of the form
  $H'/H$ for a subgroup $H'$ of $N_G(H)$. It is straightforward to
  verify that $H'$ satisfies the stated requirements.
\end{proof}
The statement about `chains of $p$-subgroups' follows immediately from
this result. Note that Cauchy's theorem and
\cref{prop:groups2:p-subgroup-chain-step} provide a new proof of
\cref{prop:groups2:sylow-i-strong} and hence of the first Sylow
theorem.

\subsection{Sylow III}
\label{subsec:groups2:sylow-iii}
The third (and last) Sylow theorem gives a good handle on the number
of $p$-Sylow subgroups of a given finite group $G$. This is especially
useful in establishing the existence of normal subgroups of $G$: all
$p$-Sylow subgroups of a group are conjugates of each other (by
the second Sylow theorem), so a $p$-Sylow subgroup is normal if and
only if $G$ has \emph{only one} $p$-Sylow
subgroup.\footnote{For the right-to-left implication in this comment,
  see the third point in \Cref{exer:groups2:characteristic-subgrp-basics}.}
\begin{thm}{Third Sylow theorem}{thm:groups2:sylow-iii}
  Let $p$ be a prime integer, and let $G$ be a finite group of order
  $|G| = p^r m$. Assume that $p$ does not divide $m$. Then the number
  of $p$-Sylow subgroups of $G$ divides $m$ and is congruent to $1$
  modulo $p$.
\end{thm}
\begin{proof}
  Let $N_p$ denote the number of $p$-Sylow subgroups of $G$.

  By \cref{thm:groups2:sylow-ii}, the $p$-Sylow subgroups of $G$ are
  the conjugates of any given $p$-Sylow subgroup $P$. By
  \cref{lem:groups2:conjugate-subgroups-count}, $N_p$ is the index of
  the normalizer $N_G(P)$ of $P$; thus
  (\cref{coro:groups2:conjugate-subgroups-divides-index}) it divides
  the index $m$ of $P$. In fact,
  \[
    m = [G : P] = [G : N_G(P)] \cdot [N_G(P) : P] = N_p \cdot [N_G(P)
    : P].
  \]
  Now, by \cref{lem:groups2:normalizer-index-mod-p} we have
  \[
    m = [G : P] \equiv [N_G(P) : P] \bmod p;
  \]
  multiplying by $N_p$, we get
  \[
    mN_p \equiv m \bmod p.
  \]
  Since $m \not\equiv 0 \bmod p$ and $p$ is prime, this implies
  \[
    N_p \equiv 1 \bmod p,
  \]
  as needed.
\end{proof}
Of course there are other ways to prove \cref{thm:groups2:sylow-iii}:
see for example \Cref{exer:groups2:alt-proof-third-sylow-thm}.

\subsection{Applications}
\label{subsec:groups2:applications}
Consequences stemming from the group actions we have encountered, and
especially the Sylow theorems, may be applied to establish exquisitely
precise facts about individual groups as well as whole classes of
groups; this is often based on some simple but clever numerology.

The following examples are exceedingly simple-minded but will
hopefully convey the flavor of what can be done with the tools we have
built in the previous two sections. More examples may be found among
the exercises at the end of this section.
\subsubsection*{More nonsimple groups}

\begin{fact}{}{fact:groups2:nonsimple-mpr}
  Let $G$ be a group of order $mp^r$, where $p$ is a prime integer and
  $1 < m < p$. Then $G$ is not simple. (Cf.\
  \Cref{exam:groups2:mp-not-simple}.)
\end{fact}
\begin{proof}
  By the third Sylow theorem, the number $N_p$ of $p$-Sylow subgroups
  divides $m$ and is of the form $1 + kp$. Since $m < p$, this forces
  $k = 0$, $N_p = 1$. Therefore $G$ has a normal subgroup of order
  $p^r$; hence it is not simple.
\end{proof}
Of course the same argument gives the same conclusion for every group
of order $mp^r$, where $\gcd(m, p) = 1$ and the only divisor $d$ of $m$
such that $d \equiv 1 \bmod p$ is $d = 1$.
\begin{exam}{}{}
  There are no simple groups of order 2002.

  Indeed,\footnote{It is safe to guess that this statement has been
    assigned on hundreds of algebra tests across the world in the year
    2002.}
  \[
    2002 = 2 \cdot 7 \cdot 11 \cdot 13;
  \]
  the divisors of $2 \cdot 7 \cdot 13$ are
  \[
    1, 2, 7, 13, 14, 26, 91, 182\,:
  \]
  of these, only $1$ is congruent to $1 \bmod 11$. Thus there is a
  normal subgroup of order 11 in every group of order 2002.
\end{exam}
The reader should not expect the third Sylow theorem to always yield
its fruits so readily, however.
\begin{exam}{}{}
  There are no simple groups of order 12.

  Note that $3 \equiv 1 \bmod 2$ and $4 \equiv 1 \bmod 3$: thus the
  argument used above does not guarantee the existence of either a
  normal 2-Sylow subgroup or a normal 3-Sylow subgroup.

  However, suppose that there is more than one 3-Sylow subgroup. Then
  there must be 4, by the third Sylow theorem. Since any two such
  subgroups must intersect in the identity, this accounts for exactly
  8 elements of order 3.  Excluding these leaves us with the identity
  and 3 elements of order 2 or 4; that is just enough room to fit one
  2-Sylow subgroup. This subgroup will then have to be normal.

  Thus, either there is a 3-Sylow normal subgroup or there is a
  2-Sylow normal subgroup---either way, the group is not simple.
\end{exam}
Even this more refined counting will often fail, and one has to dig
deeper.

\begin{exam}{}{exam:groups2:no-simple-order-24}
  There are no simple groups of order 24.

  Indeed, let $G$ be a group of order 24, and consider its 2-Sylow
  subgroups; by the third Sylow theorem, there are either 1 or 3 such
  subgroups. If there is 1, the 2-Sylow subgroup is normal and $G$ is
  not simple. Otherwise, $G$ acts (nontrivially) by conjugation on
  this set of three 2-Sylow subgroups; this action gives a nontrivial
  homomorphism $G \to S_3$, whose kernel is a proper, nontrivial
  normal subgroup of $G$---thus again $G$ is not simple.
\end{exam}
The reader should practice by selecting a random number $n$ and trying
to say as much as he/she can, in general, about groups of order
$n$. Beware: such problems are a common feature of qualifying exams.

\subsubsection*{Groups of order $pq$, $p < q$ prime}

\begin{fact}{}{fact:groups2:order-pq-cyclic}
  Assume $p < q$ are prime integers and $q \not\equiv 1 \bmod p$. Let
  $G$ be a group of order $pq$. Then $G$ is cyclic.
\end{fact}
\begin{proof}
  By the third Sylow theorem, $G$ has a unique (hence normal) subgroup
  $H$ of order $p$. Indeed, the number $N_p$ of $p$-Sylow subgroups
  must divide $q$, and $q$ is prime, so $N_p = 1$ or $q$. Necessarily
  $N_p \equiv 1 \bmod p$, and $q \not\equiv 1 \bmod p$ by hypothesis;
  therefore $N_p = 1$.  Since $H$ is normal, conjugation gives an
  action of $G$ on $H$, hence (by
  \cref{exer:groups2:conj-action-kernel-is-centralizer}) a
  homomorphism $\gamma : G \to \Aut(H)$.  Now $H$ is cyclic of order
  $p$, so $|\Aut(H)| = p - 1$ (\Cref{exer:groups:aut-cn-euler-phi});
  the order of $\gamma(G)$ must divide both $pq$ and $p - 1$, and it
  follows that $\gamma$ is the trivial map.  Therefore, conjugation is
  trivial on $H$: that is, $H \subseteq Z(G)$.
  \cref{lem:groups2:quotient-center-cyclic} implies that $G$ is
  abelian.

  Finally, an abelian group of order $pq$, with $p < q$ primes, is
  necessarily cyclic: indeed it must contain elements $g$, $h$ of
  order $p$, $q$, respectively (for example by Cauchy's theorem), and
  then $|gh| = pq$ by \Cref{exer:groups:order-product-coprime}.
\end{proof}
For example, this statement `classifies' all groups of order 15, 33,
35, 51, $\dots$: such groups are necessarily cyclic.

The argument given in the proof is rather `high-brow', as it involves
the automorphism group of $H$; that is precisely why I gave it. For
low-brow alternatives, see
\Cref{exer:groups2:alt-proof-ord-pq-cyclic-via-sylow-count} or
\Cref{rem:groups2:order-pq-cyclic-alternative-argument}.

The condition $q \not\equiv 1 \bmod p$ in
\cref{fact:groups2:order-pq-cyclic} is clearly necessary: indeed,
$|S_3| = 2 \cdot 3$ is the product of two distinct primes, and yet
$S_3$ is not cyclic. The argument given in the proof shows that if
$|G| = pq$, with $p < q$ prime, and $G$ has a normal subgroup of order
$p$, then $G$ is cyclic. If $q \equiv 1 \bmod p$, it can be shown that
there is in fact a unique noncommutative group of order $pq$ up to
isomorphism: the reader will work this out after learning about
semidirect products (\Cref{exer:groups2:classify-grps-ord-pq}). But we
are in fact already in the position of obtaining rather sophisticated
information about this group, even without knowing its construction in
general (\Cref{exer:groups2:noncomm-grp-ord-pq-structure}).

For fun, let's tackle the case in which $p = 2$.
\begin{fact}{}{fact:groups2:order-2q-dihedral}
  Let $q$ be an odd prime, and let $G$ be a noncommutative group of
  order $2q$.  Then $G \cong D_{2q}$, the dihedral group.
\end{fact}
\begin{proof}
  By Cauchy's theorem, $\exists\, y \in G$ such that $y$ has order
  $q$. By the third Sylow theorem, $\genalg{y}$ is the unique subgroup
  of order $q$ in $G$ (and is therefore normal). Since $G$ is not
  commutative and in particular it is not cyclic, it has no elements
  of order $2q$; therefore, every element in the complement of
  $\genalg{y}$ has order 2; let $x$ be any such element.  The
  conjugate $xyx^{-1}$ of $y$ by $x$ is an element of order $q$, so
  $xyx^{-1} \in \genalg{y}$. Thus, $xyx^{-1} = y^r$ for some $r$
  between 0 and $q - 1$.  Now observe that
  \[
    (y^r)^r = (xyx^{-1})^r = xy^r x^{-1} = x^2 y(x^{-1})^2 = y
  \]
  since $|x| = 2$. Therefore, $y^{r^2 - 1} = e$, which implies
  \[
    q \mid (r^2 - 1) = (r-1)(r+1)
  \]
  by \Cref{coro:groups:order-multiple}. Since $q$ is prime, this says
  that $q \mid (r-1)$ or $q \mid (r+1)$; since $0 \leq r \leq q-1$, it
  follows that $r = 1$ or $r = q - 1$.  If $r = 1$, then
  $xyx^{-1} = y$; that is, $xy = yx$. But then the order of $xy$ is
  $2q$ (by \Cref{exer:groups:order-product-coprime}), and $G$ is
  cyclic, a contradiction.

  Therefore $r = q - 1$, and we have established the relations
  \[
    \begin{cases}
      x^2 = e, \\
      y^q = e, \\
      yx = xy^{q-1}.
    \end{cases}
  \]
  These are the relations satisfied by generators $x$, $y$ of
  $D_{2q}$, as the reader hopefully verified in
  \Cref{exer:groups:dihedral-generators}; the statement follows.
\end{proof}
\cref{fact:groups2:order-2q-dihedral} yields a classification of
groups of order $2q$, for $q$ an odd prime: such a group must be
either abelian (and hence cyclic, by the usual considerations) or
isomorphic to a dihedral group.  For $q = 3$, we recover the result of
\cref{exer:groups2:noncomm-grp-ord-6-iso-s3}: every noncommutative
group of order 6 is isomorphic to $D_6 \cong S_3$.

\subsection*{Exercises}

\begin{exer}{}{exer:groups2:prove-cauchy-cyclic-subgrp-count}
  \exerused{} Prove \cref{fact:groups2:cauchy-cyclic-count}. [\Cref{subsec:groups2:cauchy-s-theorem}]
\end{exer}
\begin{exer}{}{exer:groups2:characteristic-subgrp-basics}
  \exerused{} Let $G$ be a group. A subgroup $H$ of $G$ is
  \emph{characteristic} if $\varphi(H) \subseteq H$ for every
  automorphism $\varphi$ of $G$.
  \begin{itemize}
  \item Prove that characteristic subgroups are normal.
  \item Let $H \subseteq K \subseteq G$, with $H$ characteristic in
    $K$ and $K$ normal in $G$. Prove that $H$ is normal in $G$.
  \item Let $G$, $K$ be groups, and assume that $G$ contains a single
    subgroup $H$ isomorphic to $K$. Prove that $H$ is normal in $G$.
  \item Let $K$ be a normal subgroup of a finite group $G$, and assume
    that $|K|$ and $|G/K|$ are relatively prime. Prove that $K$ is
    characteristic in $G$.
  \end{itemize} [\Cref{subsec:groups2:cauchy-s-theorem}, \Cref{subsec:groups2:sylow-iii},
  \ref{exer:groups2:sylow-normal-implies-char-self-normalizing},
  \Cref{subsec:groups2:the-commutator-subgroup-derived-series-and-solvability}]
\end{exer}
\begin{exer}{}{exer:groups2:abel-simple-iff-cyclic-prime}
  Prove that a nonzero abelian group $G$ is simple if and only if
  $G \cong \Z/p\Z$ for some positive prime integer $p$.
\end{exer}
\begin{exer}{}{exer:groups2:simple-iff-only-triv-hom-images}
  \exerused{} Prove that a nontrivial finite group $G$ is simple if
  and only if its only homomorphic images (i.e., groups $G'$ such that
  there is an onto homomorphism $G \to G'$) are the trivial group and
  $G$ itself (up to isomorphism). (Note: One implication is true for
  arbitrary groups, while $\Z[\frac{1}{2}]/\Z$ is a non-simple group
  which is isomorphic to its nontrivial homomorphic images, as the
  reader may enjoy
  verifying.) [\Cref{subsec:groups2:composition-factors-schreier-s-theorem}]
\end{exer}
\begin{exer}{}{exer:groups2:simple-grp-hom-injective}
  Let $G$ be a simple group, and assume $\varphi : G \to G'$ is a
  nontrivial group homomorphism. Prove that $\varphi$ is injective.
\end{exer}
\begin{exer}{}{exer:groups2:no-simple-p-grp-ord-geq-p2}
  Prove that there are no simple groups of order 4, 8, 9, 16, 25, 27,
  32, or 49. In fact, prove that no $p$-group of order $\geq p^2$ is
  simple.
\end{exer}
\begin{exer}{}{exer:groups2:no-simple-grp-composite-ord-list}
  Prove that there are no simple groups of order 6, 10, 14, 15, 20,
  21, 22, 26, 28, 33, 34, 35, 38, 39, 42, 44, 46, 51, 52, 55, 57, or
  58. (Hint: \Cref{exam:groups2:mp-not-simple}.)
\end{exer}
\begin{exer}{}{exer:groups2:intersect-p-sylow-normal-maximal}
  Let $G$ be a finite group, $p$ a prime integer, and let $N$ be the
  intersection of the $p$-Sylow subgroups of $G$. Prove that $N$ is a
  normal $p$-subgroup of $G$ and that every normal $p$-subgroup of $G$
  is contained in $N$. (In other words, $G/N$ is final with respect to
  the property of being a homomorphic image of $G$ of order
  $|G|/p^\alpha$ for some $\alpha$.)
\end{exer}
\begin{exer}{}{exer:groups2:p-subgrp-in-normalizer-subset-sylow}
  $\lnot$ Let $P$ be a $p$-Sylow subgroup of a finite group $G$, and
  let $H \subseteq G$ be a $p$-subgroup. Assume $H \subseteq
  N_G(P)$. Prove that $H \subseteq P$. (Hint: $P$ is normal in
  $N_G(P)$, so $PH$ is a subgroup of $N_G(P)$ by
  \Cref{prop:groups:second-isomorphism-theorem}, and
  $|PH/P| = |H/(P \cap H)|$. Show that this implies that $PH$ is a
  $p$-group, and hence $PH = P$ since $P$ is a maximal $p$-subgroup of
  $G$. Deduce that $H \subseteq P$.)
  [\ref{exer:groups2:sylow-conj-action-p-unique-fixed-pt}]
\end{exer}
\begin{exer}{}{exer:groups2:sylow-conj-action-p-unique-fixed-pt}
  $\lnot$ Let $P$ be a $p$-Sylow subgroup of a finite group $G$, and
  act with $P$ by conjugation on the set of $p$-Sylow subgroups of
  $G$. Show that $P$ is the unique fixed point of this action. (Hint:
  Use \Cref{exer:groups2:p-subgrp-in-normalizer-subset-sylow}.)
  [\ref{exer:groups2:alt-proof-third-sylow-thm}]
\end{exer}
\begin{exer}{}{exer:groups2:alt-proof-third-sylow-thm}
  \exerused{} Use the second Sylow theorem,
  \cref{coro:groups2:conjugate-subgroups-divides-index}, and
  \Cref{exer:groups2:sylow-conj-action-p-unique-fixed-pt} to paste
  together an alternative proof of the third Sylow theorem. [\Cref{subsec:groups2:sylow-iii}]
\end{exer}
\begin{exer}{}{exer:groups2:idx-containing-normalizer-cong-1-mod-p}
  Let $P$ be a $p$-Sylow subgroup of a finite group $G$, and let
  $H \subseteq G$ be a subgroup containing the normalizer
  $N_G(P)$. Prove that $[G : H] \equiv 1 \bmod p$.
\end{exer}
\begin{exer}{}{exer:groups2:sylow-normal-implies-char-self-normalizing}
  $\lnot$ Let $P$ be a $p$-Sylow subgroup of a finite group $G$.
  \begin{itemize}
  \item Prove that if $P$ is normal in $G$, then it is in fact
    characteristic in $G$ (cf.\
    \cref{exer:groups2:characteristic-subgrp-basics}).
  \item Let $H \subseteq G$ be a subgroup containing the Sylow
    subgroup $P$.  Assume $P$ is normal in $H$ and $H$ is normal in
    $G$. Prove that $P$ is normal in $G$.
  \item Prove that $N_G(N_G(P)) = N_G(P)$.
  \end{itemize}
  [\ref{exer:groups2:nilp-grp-proper-subgrp-lt-normalizer}]
\end{exer}
\begin{exer}{}{exer:groups2:no-simple-grp-ord-18-40-45-50-54}
  Prove that there are no simple groups of order 18, 40, 45, 50, or
  54.
\end{exer}
\begin{exer}{}{exer:groups2:classify-grps-ord-leq-15}
  Classify all groups of order $n \leq 15$, $n \neq 8, 12$: that is,
  produce a list of nonisomorphic groups such that every group of
  order $n \neq 8, 12$, $n \leq 15$ is isomorphic to one group in the
  list.
\end{exer}
\begin{exer}{}{exer:groups2:noncomm-grp-ord-8-is-d8-or-q8}
  $\lnot$ Let $G$ be a noncommutative group of order 8.
  \begin{itemize}
  \item Prove that $G$ contains elements of order 4 and no elements of
    order 8.
  \item Let $y$ be an element of order 4. Prove that $G$ is generated
    by $y$ and by an element $x \notin \genalg{y}$, such that
    $x^2 = e$ or $x^2 = y^2$.
  \item In either case, $G = \{e, y, y^2, y^3, x, yx, y^2x,
    y^3x\}$. Prove that the multiplication table of $G$ is determined
    by whether $x^2 = e$ or $x^2 = y^2$, and by the value of $xy$.
  \item Prove that necessarily $xy = y^3 x$. (Hint: To eliminate
    $xy = y^2 x$, multiply on the right by $y$.)
  \item Prove that $G \cong D_8$ or $G \cong Q_8$.
  \end{itemize} [\ref{exer:groups2:complete-classification-ord-8},
  \ref{exer:fields:nested-radical-galois-examples}]
\end{exer}
\begin{exer}{}{exer:groups2:division-ring-ord-64-comm}
  $\lnot$ Let $R$ be a division ring
  (\Cref{defn:rings:division-ring}), and assume $|R| = 64$. Prove that
  $R$ is necessarily commutative (hence, a field), as follows:
  \begin{itemize}
  \item The group of units of $R$ has order 63. Prove it has a
    commutative subgroup $G$ of order 9. (Sylow.)
  \item Prove that $R$ is the only sub-division ring of $R$ containing
    $G$.
  \item Prove that the set of elements of $R$ commuting with every
    element of $G$ is a sub-division ring of $R$ containing $G$. (Cf.\
    \Cref{exer:rings:centralizer-is-subring}.)
  \item Conclude that $G$ is contained in the center of $R$. Recall
    that the center of $R$ is a sub-division ring of $R$ (cf.\
    \Cref{exer:rings:center-is-subring}), and conclude that $R$ is
    commutative.
  \end{itemize}
  Like \Cref{exer:rings:division-ring-p2-commutative}, this is a
  particular case of a theorem of Wedderburn, according to which every
  finite division ring is a
  field. [\ref{exer:fields:wedderburn-thm-birkhoff-vandiver-proof}]
\end{exer}
\begin{exer}{}{exer:groups2:alt-proof-ord-pq-cyclic-via-sylow-count}
  \exerused{} Give an alternative proof of
  \cref{fact:groups2:order-pq-cyclic} as follows: use the third Sylow
  theorem to count the number of elements of order $p$ and $q$ in $G$;
  use this to show that there are elements in $G$ of order neither 1
  nor $p$ nor $q$; deduce that $G$ is cyclic. [\Cref{subsec:groups2:applications}]
\end{exer}
\begin{exer}{}{exer:groups2:noncomm-grp-ord-pq-structure}
  \exerused{} Let $G$ be a noncommutative group of order $pq$, where
  $p < q$ are primes.
  \begin{itemize}
  \item Show that $q \equiv 1 \bmod p$.
  \item Show that the center of $G$ is trivial.
  \item Draw the lattice of subgroups of $G$.
  \item Find the number of elements of each possible order in $G$.
  \item Find the number and size of the conjugacy classes in $G$.
  \end{itemize} [\Cref{subsec:groups2:applications}]
\end{exer}
\begin{exer}{}{exer:groups2:ord-7-elts-in-simple-grp-168}
  How many elements of order 7 are there in a simple group of order
  168?
\end{exer}
\begin{exer}{}{exer:groups2:grp-ord-pqr-not-simple}
  Let $p < q < r$ be prime integers, and let $G$ be a group of order
  $pqr$.  Prove that $G$ is not simple.
\end{exer}

\begin{exer}{}{exer:groups2:unique-divisor-cong-1-implies-not-simple}
  Let $G$ be a finite noncommutative group, $n = |G|$, and $p$ be a
  prime divisor of $n$.  Assume that the only divisor of $n$ that is
  congruent to $1$ modulo $p$ is $1$. Prove that $G$ is not simple.
\end{exer}

\begin{exer}{}{exer:groups2:simple-grp-ord-divides-np-factorial}
  $\lnot$ Let $N_p$ denote the number of $p$-Sylow subgroups of a
  group $G$.  Prove that if a noncommutative group $G$ is simple, then
  $|G|$ divides $N_p!$
  for all primes $p$ in the factorization of $|G|$. More generally,
  prove that if $G$ is simple and $H$ is a subgroup of $G$ of index
  $N > 1$, then $|G|$ divides $N!$. (Hint:
  \Cref{exer:groups:index-n-normal-subgroup-bound}.) This problem
  capitalizes on the idea behind
  \Cref{exam:groups2:no-simple-order-24}.
  [\ref{exer:groups2:simple-grp-ord-60-subgrp-idx-5}]
\end{exer}

\begin{exer}{}{exer:groups2:no-noncomm-simple-below-60}
  \exerused{} Prove that there are no noncommutative simple groups of
  order less than 60. If you have sufficient stamina, prove that the
  next possible order for a noncommutative simple group is 168. (Don't
  feel too bad if you have to cheat and look up a few particularly
  troublesome orders $> 60$.)
  [\Cref{subsec:groups2:conjugacy-an-simplicity}]
\end{exer}

\begin{exer}{}{exer:groups2:simple-grp-ord-60-subgrp-idx-5}
  $\lnot$ Assume that $G$ is a simple group of order 60.
  \begin{itemize}
  \item Use Sylow's theorems and simple numerology to prove that $G$
    has either five or fifteen 2-Sylow subgroups, accounting for
    fifteen elements of order 2 or
    4. (\cref{exer:groups2:simple-grp-ord-divides-np-factorial} will
    likely be helpful.)
  \item Show that in fact $G$ has \emph{exactly five} 2-Sylow
    subgroups.
  \end{itemize}
  Conclude that every simple group\footnote{The reader will prove
    later (\Cref{exer:groups2:a5-unique-simple-grp-ord-60}) that there
    is in fact only one simple group of order 60 up to isomorphism and
    that this group contains exactly five 2-Sylow subgroups. The
    result obtained here will be needed to establish this fact.} of
  order 60 contains a subgroup of index
  5. [\ref{exer:groups2:a5-unique-simple-grp-ord-60}]
\end{exer}

\section{Composition series and solvability}
\label{sec:groups2:composition-series-and-solvability}


I have claimed that \emph{simple} groups (in the sense of
\cref{defn:groups2:simple-group}) are the `basic constituents' of all
finite groups. Among other things, the material in this section will
(partially) justify this claim.

\subsection{The Jordan-H\"older theorem}
\label{subsec:groups2:the-jordan-h-older-theorem}
A \emph{series} of subgroups $G_i$ of a group $G$ is a decreasing
sequence of subgroups starting from $G$:
\[
  G = G_0 \supsetneq G_1 \supsetneq G_2 \supsetneq \dots
\]
The \emph{length} of a series is the number of strict inclusions.

A series is \emph{normal} if $G_{i+1}$ is normal in $G_i$ for all
$i$. We will be interested in the maximal length of a normal series in
$G$; if finite, I will denote this number\footnote{There does not
  appear to be a standard notation for this concept.} by
$\ell(G)$. The number $\ell(G)$ is a measure of how far $G$ is from
being simple. Indeed, $\ell(G) = 0$ if and only if $G$ is trivial, and
$\ell(G) = 1$ if and only if $G$ is simple: for a simple group, the
only maximal normal series is
\[
  G \supsetneq \{e\}.
\]
\begin{defn}{}{defn:groups2:composition-series}
  A \emph{composition series} for $G$ is a normal series
  \[
    G = G_0 \supsetneq G_1 \supsetneq G_2 \supsetneq \dots \supsetneq
    G_n = \{e\}
  \]
  such that the successive quotients $G_i/G_{i+1}$ are simple.
\end{defn}
It is clear (by induction on the order) that finite groups have
composition series, while infinite groups do not necessarily have one
(\Cref{exer:groups2:fin-grp-has-comp-series-z-does-not}). It
is also clear that if a normal series has maximal length $\ell(G)$,
then it is a composition series. What is not clear is that the
converse holds: conceivably, there could exist composition series of
different lengths (the longest ones having length $\ell(G)$). For
example, why can't there be a finite group $G$ with $\ell(G) = 3$ and
two different composition series
\[
  G \supsetneq G_1 \supsetneq G_2 \supsetneq \{e\}
\]
and
\[
  G \supsetneq G'_1 \supsetneq \{e\}
\]
(that is: a finite group $G$ with $\ell(G) = 3$ and a simple normal
subgroup $G'_1$ such that $G/G'_1$ is simple)?

Part of the content of the Jordan-H\"older theorem is that (luckily)
this cannot happen. In fact, the theorem is much more precise: not
only do all composition series have the same length, but they also
have the same quotients (appearing, however, in possibly different
orders).
\begin{thm}{Jordan-H\"older}{thm:groups2:jordan-holder}
  Let $G$ be a group, and let
  \begin{align*}
    G &= G_0 \supsetneq G_1 \supsetneq G_2 \supsetneq \dots
        \supsetneq G_n = \{e\},\\
    G &= G'_0 \supsetneq G'_1 \supsetneq G'_2 \supsetneq \dots
        \supsetneq G'_m = \{e\}
  \end{align*}
  be two composition series for $G$. Then $m = n$, and the lists of
  quotient groups $H_i = G_i/G_{i+1}$, $H'_i = G'_i/G'_{i+1}$ agree
  (up to isomorphism) after a permutation of the indices.
\end{thm}
\begin{proof}
  Let
  \[
    \tag{*} G = G_0 \supsetneq G_1 \supsetneq G_2 \supsetneq \dots
    \supsetneq G_n = \{e\}
  \]
  be a composition series. Argue by induction on $n$: if $n = 0$, then
  $G$ is trivial, and there is nothing to prove. Assume $n > 0$, and
  let
  \[
    \tag{**} G = G'_0 \supsetneq G'_1 \supsetneq G'_2 \supsetneq \dots
    \supsetneq G'_m = \{e\}
  \]
  be another composition series for $G$. If $G_1 = G'_1$, then the
  result follows from the induction hypothesis, since $G_1$ has a
  composition series of length $n - 1 < n$.  We may then assume
  $G_1 \neq G'_1$. Note that $G_1 G'_1 = G$: indeed, $G_1 G'_1$ is
  normal in $G$
  (\Cref{exer:groups2:product-of-normal-subgrps-normal}), and
  $G_1 \subsetneq G_1 G'_1$; but there are no proper normal subgroups
  between $G_1$ and $G$ since $G/G_1$ is simple.  Let
  $K = G_1 \cap G'_1$. The distinct subgroups $G_i \cap K$ determine a
  composition series
  \[
    K \supsetneq K_1 \supsetneq K_2 \supsetneq \dots \supsetneq K_r =
    \{e\}
  \]
  of $K$: this is not difficult to see, and will be verified more
  formally in the proof of \cref{prop:groups2:comp-series-normal}. By
  \Cref{prop:groups:second-isomorphism-theorem} (the ``second
  isomorphism theorem''),
  \[
    \frac{G_1}{K} = \frac{G_1}{G_1 \cap G'_1} \cong \frac{G_1
      G'_1}{G'_1} = \frac{G}{G'_1} \quad\text{and}\quad \frac{G'_1}{K}
    \cong \frac{G}{G_1}
  \]
  are simple. Therefore, we have two new composition series for $G$:
  \[
    \begin{array}{ccccccccc}
      G & \supsetneq & G_1 & \supsetneq & K & \supsetneq & K_1
      & \supsetneq \cdots \supsetneq & \{e\}\\
      \| & & & & \| & & \| & & \|\\
      G & \supsetneq & G'_1 & \supsetneq & K & \supsetneq & K_1
      & \supsetneq \cdots \supsetneq & \{e\}
    \end{array}
  \]
  which only differ at the first step. These two series trivially have
  the same length and the same quotients (the first two quotients get
  switched from one series to the other).  Now I claim that the first
  of these two series has the same length and quotients as the series
  $(*)$. Indeed,
  \[
    G_1 \supsetneq K \supsetneq K_1 \supsetneq K_2 \supsetneq \dots
    \supsetneq K_r = \{e\}
  \]
  is a composition series for $G_1$: by the induction hypothesis, it
  must have the same length and quotients as the composition series
  \[
    G_1 \supsetneq G_2 \supsetneq \dots \supsetneq G_n = \{e\};
  \]
  verifying my claim (and note that, in particular, $r = n - 2$).  By
  the same token, applying the induction hypothesis to the series (of
  length $n - 1$)
  \[
    G'_1 \supsetneq K \supsetneq K_1 \supsetneq K_2 \supsetneq \dots
    \supsetneq K_{n-2} = \{e\},
  \]
  shows that the second series has the same length and quotients as
  $(**)$, and the statement follows.
\end{proof}

\subsection{Composition factors; Schreier's theorem}
\label{subsec:groups2:composition-factors-schreier-s-theorem}
Two normal series are \emph{equivalent} if they have the same length
and the same quotients (up to order). The Jordan-H\"older theorem
shows that any two maximal finite series of a group are
equivalent. That is, the (isomorphism classes of the) quotients of a
composition series depend only on the group, not on the chosen
series. These are the \emph{composition factors} of the group. They
form a
multiset\footnote{See~\Cref{subsec:prelims:examples-multisets-indexed-sets}
  for a reminder on multisets: they are sets of elements counted with
  multiplicity.} of simple groups: the `basic constituents' of our
loose comment back in~\Cref{sec:groups2:the-sylow-theorems}.  It is
clear that two isomorphic groups must have the same composition
factors.  Unfortunately, it is not possible to reconstruct a group
from its composition factors alone
(\Cref{exer:groups2:nonisom-grps-same-comp-factors}). One has
to take into account the way the simple groups are `glued' together;
we will come back to this point
in~\Cref{subsec:groups2:exact-sequences-of-groups-extension-problem}.

The intuition that the composition factors of a group are its basic
constituents is reinforced by the following fact: if $G$ is a group
with a composition series, then the composition factors of every
normal subgroup $N$ of $G$ are composition factors of $G$ and the
remaining ones are the composition factors of the quotient $G/N$. For
example, the composition factors of $\Z/4\Z$ form the multiset
consisting of two copies of $\Z/2\Z$.
\begin{exam}{}{exam:groups2:z6-composition-series}
  Let $G = \Z/6\Z = \{[0],[1],[2],[3],[4],[5]\}$. Then
  \[
    \{[0],[1],[2],[3],[4],[5]\} \supsetneq \{[0],[3]\} \supsetneq
    \{[0]\}
  \]
  is a composition series for $G$; the quotients are $\Z/3\Z$,
  $\Z/2\Z$, respectively. The (normal) subgroup $N = \{[0],[2],[4]\}$
  `turns off' the second factor: indeed, intersecting the series with
  $N$ gives
  \[
    \{[0],[2],[4]\} \supsetneq \{[0]\} = \{[0]\},
  \]
  a series with composition factor $\Z/3\Z$. On the other hand,
  `mod-ing out by $N$' turns off the first factor: keeping in mind
  $[3]+N = [1]+N$, etc., we find
  \[
    \{[0]+N,[1]+N\} \supsetneq \{[0]+N\},
  \]
  a series with lone composition factor $\Z/2\Z$.
\end{exam}

This phenomenon holds in complete generality:
\begin{prop}{}{prop:groups2:comp-series-normal}
  Let $G$ be a group, and let $N$ be a normal subgroup of $G$. Then
  $G$ has a composition series if and only if both $N$ and $G/N$ have
  composition series.  Further, if this is the case, then
  \[
    \ell(G) = \ell(N) + \ell(G/N),
  \]
  and the composition factors of $G$ consist of the collection of
  composition factors of $N$ and of $G/N$.
\end{prop}
\begin{proof}
  If $G/N$ has a composition series, the subgroups appearing in it
  correspond to subgroups of $G$ containing $N$, with isomorphic
  quotients, by \Cref{prop:groups:third-isomorphism-theorem} (the
  ``third isomorphism theorem''). Thus, if both $G/N$ and $N$ have
  composition series, juxtaposing them produces a composition series
  for $G$, with the stated consequence on composition factors.  The
  converse is a little trickier. Assume that $G$ has a composition
  series
  \[
    G = G_0 \supsetneq G_1 \supsetneq G_2 \supsetneq \dots \supsetneq
    G_n = \{e\}
  \]
  and that $N$ is a normal subgroup of $G$. Intersecting the series
  with $N$ gives a sequence of subgroups of the latter:
  \[
    N = G \cap N \supseteq G_1 \cap N \supseteq \dots \supseteq \{e\}
    \cap N = \{e\}
  \]
  such that $G_{i+1} \cap N$ is normal in $G_i \cap N$, for all $i$. I
  claim that this becomes a composition series for $N$ once
  repetitions are eliminated. Indeed, this follows once we establish
  that
  \[
    \frac{G_i \cap N}{G_{i+1} \cap N}
  \]
  is either trivial (so that $G_{i+1} \cap N = G_i \cap N$, and the
  corresponding inclusion may be omitted) or isomorphic to
  $G_i/G_{i+1}$ (hence simple, and one of the composition factors of
  $G$). To see this, consider the homomorphism
  \[
    G_i \cap N \hookrightarrow G_i \twoheadrightarrow
    \frac{G_i}{G_{i+1}}:
  \]
  the kernel is clearly $G_{i+1} \cap N$; therefore (by the first
  isomorphism theorem) we have an injective homomorphism
  \[
    \frac{G_i \cap N}{G_{i+1} \cap N} \hookrightarrow
    \frac{G_i}{G_{i+1}}
  \]
  identifying $(G_i \cap N)/(G_{i+1} \cap N)$ with a subgroup of
  $G_i/G_{i+1}$.  Now, this subgroup is normal (because $N$ is normal
  in $G$) and $G_i/G_{i+1}$ is simple; our claim follows.  As for
  $G/N$, obtain a sequence of subgroups from a composition series for
  $G$:
  \[
    \frac{G}{N} \supseteq \frac{G_1 N}{N} \supseteq \frac{G_2 N}{N}
    \supseteq \cdots \supseteq \frac{\{e_G\}N}{N} = \{e_{G/N}\},
  \]
  such that $(G_{i+1}N)/N$ is normal in $(G_i N)/N$. As above, we have
  to check that
  \[
    \frac{(G_i N)/N}{(G_{i+1} N)/N}
  \]
  is either trivial or isomorphic to $G_i/G_{i+1}$. By the third
  isomorphism theorem, this quotient is isomorphic to
  $(G_i N)/(G_{i+1} N)$. This time, consider the homomorphism
  \[
    G_i \hookrightarrow G_i N \twoheadrightarrow \frac{G_i N}{G_{i+1}
      N}:
  \]
  this is surjective (check!), and the subgroup $G_{i+1}$ of the
  source is sent to the identity element in the target; hence (by
  \Cref{thm:groups:quotient-universal-property}) there is an onto
  homomorphism
  \[
    \frac{G_i}{G_{i+1}} \twoheadrightarrow \frac{G_i N}{G_{i+1} N}.
  \]
  Since $G_i/G_{i+1}$ is simple, it follows that $(G_i N)/(G_{i+1} N)$
  is either trivial or isomorphic to it
  (\Cref{exer:groups2:simple-iff-only-triv-hom-images}), as
  needed.  Summarizing, we have shown that if $G$ has a composition
  series and $N$ is normal in $G$, then both $N$ and $G/N$ have
  composition series. The first part of the argument yields the
  statement on lengths and composition factors, concluding the proof.
\end{proof}
One nice consequence of the Jordan-H\"older theorem is the following
observation. A series is a \emph{refinement} of another series if all
terms of the second appear in the first.

\begin{prop}{}{prop:groups2:schreier-finite}
  Any two normal series of a finite group ending with $\{e\}$ admit
  equivalent refinements.
\end{prop}

\begin{proof}
  Refine the series to a composition series; then apply the
  Jordan-H\"older theorem.
\end{proof}
In fact, \emph{Schreier's theorem} asserts that this holds for all
groups (while the argument given here only works for groups admitting
a composition series, e.g., finite groups). Proving this in general is
reasonably straightforward, from judicious applications of the second
isomorphism theorem (cf.\
\Cref{exer:groups2:alt-proof-jordan-holder-via-schreier}).

\subsection{The commutator subgroup, derived series, and solvability}
\label{subsec:groups2:the-commutator-subgroup-derived-series-and-solvability}
It has been a while since we have encountered a universal object; here
is one. For any group $G$, consider the category whose objects are
group homomorphisms $\alpha : G \to A$ from $G$ to a
\emph{commutative} group and whose morphisms $\alpha \to \beta$ are
(as the reader should expect) commutative diagrams
\[
  \begin{tikzcd}
    & G \arrow[dl, "\alpha"'] \arrow[dr, "\beta"] & \\
    A \arrow[rr, "\varphi"'] & & B
  \end{tikzcd}
\]
where $\varphi$ is a homomorphism.  Does this category have an initial
object? That is, given a group $G$, does there exist a
\emph{commutative} group which is universal with respect to the
property of being a homomorphic image of $G$?

Yes.

Such a group may well be thought of as the closest `commutative
approximation' of the given group $G$. To verify that this universal
object exists, we introduce the following important notion. (The
diligent reader has begun exploring this territory already, in
\Cref{exer:groups:commutator-subgroup-normal}.)
\begin{defn}{}{defn:groups2:commutator-subgroup}
  Let $G$ be a group. The \emph{commutator subgroup} of $G$ is the
  subgroup generated by all elements $ghg^{-1}h^{-1}$ with
  $g, h \in G$.
\end{defn}
The element $ghg^{-1}h^{-1}$ is often denoted $[g,h]$ and is called
the \emph{commutator} of $g$ and $h$. Thus, $g, h$ commute with each
other if and only if $[g,h] = e$.

In the same notational style, the commutator subgroup of $G$ should be
denoted $[G,G]$; this is a bit heavy, and the common shorthand for it
is $G'$, which offers the possibility of iterating the notation. Thus,
$G''$ may be used to denote the commutator subgroup of the commutator
subgroup of $G$, and $G^{(i)}$ denotes the $i$-th iterate. I will
adopt this notation in this subsection for convenience, but not
elsewhere in this book (as I want to be able to `prime' any letter I
wish, for any reason).  First we record the following trivial, but
useful, remark:

\begin{lem}{}{lem:groups2:commutator-homomorphism}
  Let $\varphi : G_1 \to G_2$ be a group homomorphism. Then
  $\forall g, h \in G_1$ we have
  \[
    \varphi([g,h]) = [\varphi(g), \varphi(h)]
  \]
  and $\varphi(G'_1) \subseteq G'_2$.
\end{lem}
This simple observation makes the key properties of the commutator
subgroup essentially immediate (cf.\
\Cref{exer:groups:commutator-subgroup-normal}):

\begin{prop}{}{prop:groups2:commutator-subgroup-props}
  Let $G'$ be the commutator subgroup of $G$. Then
  \begin{itemize}
  \item $G'$ is normal in $G$;
  \item the quotient $G/G'$ is commutative;
  \item if $\alpha : G \to A$ is a homomorphism of $G$ to a
    commutative group, then $G' \subseteq \ker\alpha$;
  \item the natural projection $G \to G/G'$ is universal in the sense
    explained above.
  \end{itemize}
\end{prop}
\begin{proof}
  These are all easy consequences of
  \cref{lem:groups2:commutator-homomorphism}:

  ---By \cref{lem:groups2:commutator-homomorphism}, the commutator
  subgroup is characteristic, hence normal (cf.\
  \Cref{exer:groups2:characteristic-subgrp-basics}).

  ---By \cref{lem:groups2:commutator-homomorphism}, the commutator of
  any two cosets $gG'$, $hG'$ is the coset of the commutator $[g,h]$;
  hence it is the identity in $G/G'$. As noted above, this implies
  that $G/G'$ is commutative.

  ---Let $\alpha : G \to A$ be a homomorphism to a commutative
  group. By \cref{lem:groups2:commutator-homomorphism},
  $\alpha(G') \subseteq A' = \{e\}$: that is,
  $G' \subseteq \ker\alpha$.

  ---The universality follows from the previous point and from the
  universal property of quotients
  (\Cref{thm:groups:quotient-universal-property}).
\end{proof}
Taking successive commutators of a group produces a descending
sequence of subgroups,
\[
  G \supseteq G' \supseteq G'' \supseteq G''' \supseteq \dots,
\]
which is `normal' in the sense indicated
in~\Cref{subsec:groups2:the-jordan-h-older-theorem}.
\begin{defn}{}{defn:groups2:derived-series}
  Let $G$ be a group. The \emph{derived series} of $G$ is the sequence
  of subgroups
  \[
    G \supseteq G' \supseteq G'' \supseteq G''' \supseteq \dots.
  \]
\end{defn}
The derived series may or may not end with the identity of $G$. For
example, if $G$ is commutative, then the sequence gets there right
away:
\[
  G \supseteq G' = \{e\};
\]
however, if $G$ is simple and noncommutative, then it gets stuck at
the first step:
\[
  G = G' = G'' = \cdots
\]
(indeed, $G'$ is normal and $\neq \{e\}$ as $G$ is noncommutative; but
then $G' = G$ since $G$ is simple).
\begin{defn}{}{defn:groups2:solvable}
  A group is \emph{solvable} if its derived series terminates with the
  identity.
\end{defn}
For example, abelian groups are solvable.

The importance of this notion will be most apparent in the relatively
distant future because of a brilliant application of Galois theory
(\Cref{subsec:fields:solvability-of-polynomial-equations-by-radicals}).
But we can already appreciate it in the way it relates to the material
we just covered. A normal series is \emph{abelian}, resp.,
\emph{cyclic}, if all quotients are abelian, resp.,
cyclic\footnote{Thus, a composition series should be called `simple';
  to our knowledge, it is not.}.
\begin{prop}{}{prop:groups2:solvable-equiv}
  For a finite group $G$, the following are equivalent:
  \begin{enumerate}[label=(\roman*)]
  \item All composition factors of $G$ are cyclic.
  \item $G$ admits a cyclic series ending in $\{e\}$.
  \item $G$ admits an abelian series ending in $\{e\}$.
  \item $G$ is solvable.
  \end{enumerate}
\end{prop}
\begin{proof}
  (i) $\implies$ (ii) $\implies$ (iii) are trivial. (iii) $\implies$
  (i) is obtained by refining an abelian series to a composition
  series (keeping in mind that the simple abelian groups are cyclic
  $p$-groups).

  (iv) $\implies$ (iii) is also trivial, since the derived series is
  abelian (by the second point in
  \cref{prop:groups2:commutator-subgroup-props}).  Thus, we only have
  to prove (iii) $\implies$ (iv). For this, let
  \[
    G = G_0 \supsetneq G_1 \supsetneq G_2 \supsetneq \dots \supsetneq
    G_n = \{e\}
  \]
  be an abelian series. Then I claim that $G^{(i)} \subseteq G_i$ for
  all $i$, where $G^{(i)}$ denotes the $i$-th `iterated' commutator
  subgroup.  This can be verified by induction. For $i = 1$, $G/G_1$
  is commutative; thus $G' \subseteq G_1$, by the third point in
  \cref{prop:groups2:commutator-subgroup-props}. Assuming we know
  $G^{(i)} \subseteq G_i$, the fact that $G_i/G_{i+1}$ is abelian
  implies $G'_i \subseteq G_{i+1}$, and hence
  \[
    G^{(i+1)} = (G^{(i)})' \subseteq G'_i \subseteq G_{i+1},
  \]
  as claimed.  In particular we obtain that
  $G^{(n)} \subseteq G_n = \{e\}$: that is, the derived series
  terminates at $\{e\}$, as needed.
\end{proof}
\begin{exam}{}{exam:groups2:p-groups-are-solvable}
  All $p$-groups are solvable. Indeed, the composition factors of a
  $p$-group are simple $p$-groups (what else could they be?), hence
  cyclic.
\end{exam}
\begin{coro}{}{coro:groups2:solvable-normal}
  Let $N$ be a normal subgroup of a finite group $G$. Then $G$ is
  solvable if and only if both $N$ and $G/N$ are solvable.
\end{coro}

\begin{proof}
  This follows immediately from \cref{prop:groups2:comp-series-normal}
  and the formulation of solvability in terms of composition factors
  given in \cref{prop:groups2:solvable-equiv}.
\end{proof}
It is worth mentioning that any subgroup $H$ of a solvable group is
solvable: indeed, the commutator $H'$ of $H$ is a subgroup of the
commutator $G'$ of $G$, hence $H'' \subseteq G''$,
$H''' \subseteq G'''$, and so on.

The \emph{Feit-Thompson theorem} asserts that every finite group of
odd order is solvable. This result is many orders of magnitude beyond
the scope of this book: the original 1963 proof runs about 250 pages.

\subsection*{Exercises}

\begin{exer}{}{exer:groups2:z-normal-series-arb-length}
  Prove that $\Z$ has normal series of arbitrary lengths. (Thus,
  $\ell(\Z)$ is not finite.)
\end{exer}

\begin{exer}{}{exer:groups2:length-of-fin-cyclic-grp}
  Let $G$ be a finite \emph{cyclic} group. Compute $\ell(G)$ in terms
  of $|G|$. Generalize to finite solvable groups.
\end{exer}
\begin{exer}{}{exer:groups2:fin-grp-has-comp-series-z-does-not}
  \exerused Prove that every finite group has a composition
  series. Prove that $\Z$ does not have a composition
  series. [\Cref{subsec:groups2:the-jordan-h-older-theorem}]
\end{exer}

\begin{exer}{}{exer:groups2:nonisom-grps-same-comp-factors}
  \exerused Find an example of two nonisomorphic groups with the same
  composition
  factors. [\Cref{subsec:groups2:composition-factors-schreier-s-theorem}]
\end{exer}
\begin{exer}{}{exer:groups2:product-of-normal-subgrps-normal}
  \exerused Show that if $H$, $K$ are \emph{normal} subgroups of a
  group $G$, then $HK$ is a normal subgroup of
  $G$. [\Cref{subsec:groups2:the-jordan-h-older-theorem}]
\end{exer}

\begin{exer}{}{exer:groups2:product-has-comp-series-iff-both}
  Prove that $G_1 \times G_2$ has a composition series if and only if
  both $G_1$ and $G_2$ do, and explain how the corresponding
  composition factors are related.
\end{exer}
\begin{exer}{}{exer:groups2:alt-proof-jordan-holder-via-schreier}
  \exerused Locate and understand a proof of (the general form of)
  Schreier's theorem that does not use the Jordan-H\"older
  theorem. Then obtain an alternative proof of the Jordan-H\"older
  theorem, using
  Schreier's. [\Cref{subsec:groups2:composition-factors-schreier-s-theorem}]
\end{exer}

\begin{exer}{}{exer:groups2:prove-commutator-hom-lem}
  \exerused Prove
  \Cref{lem:groups2:commutator-homomorphism}. [\Cref{subsec:groups2:the-commutator-subgroup-derived-series-and-solvability}]
\end{exer}
\begin{exer}{}{exer:groups2:p-grp-abelian-series-construction}
  Let $G$ be a nontrivial $p$-group. Construct explicitly an abelian
  series for $G$, using the fact that the center of a nontrivial
  $p$-group is nontrivial
  (\cref{coro:groups2:p-group-nontrivial-center}). This gives an
  alternative proof of the fact that $p$-groups are solvable
  (\cref{exam:groups2:p-groups-are-solvable}).
\end{exer}
\begin{exer}{}{exer:groups2:nilpotent-grp-upper-central-series}
  $\lnot$ Let $G$ be a group. Define inductively an increasing
  sequence $Z_0 = \{e\} \subseteq Z_1 \subseteq Z_2 \subseteq \dots$
  of subgroups of $G$ as follows: for $i \geq 1$, $Z_i$ is the
  subgroup of $G$ corresponding (as in
  \Cref{prop:groups:subgroups-of-quotients-bijection}) to the center
  of $G/Z_{i-1}$.
  \begin{itemize}
  \item Prove that each $Z_i$ is normal in $G$, so that this
    definition makes sense.
  \end{itemize}
  A group is\footnote{There are many alternative characterizations for
    this notion that are equivalent to the one given here but not too
    trivially so.}  \emph{nilpotent} if $Z_m = G$ for some $m$.
  \begin{itemize}
  \item Prove that $G$ is nilpotent if and only if $G/Z(G)$ is
    nilpotent.
  \item Prove that $p$-groups are nilpotent.
  \item Prove that nilpotent groups are solvable.
  \item Find a solvable group that is not nilpotent.
  \end{itemize} [\ref{exer:groups2:nilp-grp-normal-subgrp-meets-ctr},
  \ref{exer:groups2:nilp-grp-proper-subgrp-lt-normalizer},
  \ref{exer:groups2:all-sylow-normal-implies-direct-product}]
\end{exer}
\begin{exer}{}{exer:groups2:nilp-grp-normal-subgrp-meets-ctr}
  $\lnot$ Let $H$ be a nontrivial normal subgroup of a nilpotent group
  $G$ (cf.\
  \Cref{exer:groups2:nilpotent-grp-upper-central-series}). Prove
  that $H$ intersects $Z(G)$ nontrivially. (Hint: Let $r \geq 1$ be
  the smallest index such that $\exists h \neq e$, $h \in H \cap
  Z_r$. Contemplate a well-chosen commutator $[g,h]$.) Since
  $p$-groups are nilpotent, this strengthens the result of
  \Cref{exer:groups2:nontriv-norm-subgrp-meets-ctr}.
  [\ref{exer:groups2:solvable-grp-normal-subgrp-has-comm-normal}]
\end{exer}
\begin{exer}{}{exer:groups2:nilp-grp-proper-subgrp-lt-normalizer}
  Let $H$ be a proper subgroup of a finite nilpotent group $G$ (cf.\
  \Cref{exer:groups2:nilpotent-grp-upper-central-series}). Prove
  that $H \subsetneq N_G(H)$. (Hint: $Z(G)$ is nontrivial. First
  dispose of the case in which $H$ does not contain $Z(G)$, and then
  use induction to deal with the case in which $H$ does contain
  $Z(G)$.) Deduce that every Sylow subgroup of a finite nilpotent
  group is normal\footnote{This property characterizes finite
    nilpotent groups; cf.\
    \Cref{exer:groups2:all-sylow-normal-implies-direct-product}.}. (Use
  \Cref{exer:groups2:sylow-normal-implies-char-self-normalizing}.)
\end{exer}
\begin{exer}{}{exer:groups2:derived-series-terms-characteristic}
  $\lnot$ For a group $G$, let $G^{(i)}$ denote the iterated
  commutator, as
  in~\Cref{subsec:groups2:the-commutator-subgroup-derived-series-and-solvability}. Prove
  that each $G^{(i)}$ is characteristic (hence normal) in
  $G$. [\ref{exer:groups2:solvable-grp-normal-subgrp-has-comm-normal}]
\end{exer}
\begin{exer}{}{exer:groups2:solvable-grp-normal-subgrp-has-comm-normal}
  Let $H$ be a nontrivial normal subgroup of a solvable group $G$.
  \begin{itemize}
  \item Prove that $H$ contains a nontrivial commutative subgroup that
    is normal in $G$. (Hint: Let $r$ be the largest index such that
    $K = H \cap G^{(r)}$ is nontrivial. Prove that $K$ is commutative,
    and use
    \Cref{exer:groups2:derived-series-terms-characteristic} to
    show it is normal in $G$.)
  \item Find an example showing that $H$ need not intersect the center
    of $G$ nontrivially (cf.\
    \Cref{exer:groups2:nilp-grp-normal-subgrp-meets-ctr}).
  \end{itemize}
\end{exer}
\begin{exer}{}{exer:groups2:grp-ord-p2q-solvable}
  Let $p$, $q$ be prime integers, and let $G$ be a group of order
  $p^2 q$.  Prove that $G$ is solvable. (This is a particular case of
  Burnside's theorem: for $p$, $q$ primes, every group of order
  $p^a q^b$ is solvable.)
\end{exer}
\begin{exer}{}{exer:groups2:grp-ord-lt-120-neq-60-solvable}
  \exerused Prove that every group of order $< 120$ and $\neq 60$ is
  solvable.  [\Cref{subsec:groups2:conjugacy-an-simplicity},
  \Cref{subsec:fields:solvability-of-polynomial-equations-by-radicals}]
\end{exer}

\begin{exer}{}{exer:groups2:feit-thompson-equiv-even-ord}
  Prove that the Feit-Thompson theorem is equivalent to the assertion
  that every noncommutative finite simple group has even order.
\end{exer}

\section{The symmetric group}
\label{sec:groups2:the-symmetric-group}


\subsection{Cycle notation}
\label{subsec:groups2:cycle-notation}
It is time to give a second look at symmetric groups. Recall that
$S_n$ denotes the group of permutations (i.e., automorphisms in
$\categ{Set}$) of the set
$\{1,\dots,n\}$. In~\Cref{sec:groups:examples-of-groups} we denoted
elements of $S_n$ in a straightforward but inconvenient way:
\[
  \sigma = \begin{pmatrix} 1 & 2 & 3 & 4 & 5 & 6 & 7 & 8 \\
    8 & 1 & 2 & 7 & 5 & 3 & 4 & 6 \end{pmatrix}
\]
would stand for the element in $S_8$ sending $\mathbf{1}$ to
$\mathbf{8}$, $\mathbf{2}$ to $\mathbf{1}$, etc. (Here I return to the
convention in \Cref{sec:groups:examples-of-groups} and let $S_A$ act
on $A$ on the \emph{right}.)

There is clearly ``too much'' information here (the first row should
be implicit), and at the same time it seems hard to find out anything
interesting about a permutation from this notation. For example, can
the reader say anything about the conjugates of $\sigma$ in $S_8$? For
maximal enlightenment, try to do
\Cref{exer:groups2:conj-class-size-example-s8} now, and then
try again after absorbing the material
in~\Cref{subsec:groups2:type-conjugacy-sn}. In short, we should be
able to do better.  As often is the case, thinking in terms of actions
helps. By its very definition, the group $S_n$ acts on the set
$\{1,\dots,n\}$; so does every subgroup of $S_n$. Given a permutation
$\sigma \in S_n$, consider the cyclic group $\langle\sigma\rangle$
generated by $\sigma$ and its action on $\{1,\dots,n\}$. The orbits of
this action form a partition of $\{1,\dots,n\}$; therefore, every
$\sigma \in S_n$ determines a partition of $\{1,\dots,n\}$. For
example, the element $\sigma \in S_8$ given above splits
$\{1,\dots,8\}$ into three orbits:
\[
  \{1,2,3,6,8\}, \quad \{4,7\}, \quad \{5\}.
\]
The action of $\langle\sigma\rangle$ is transitive on each orbit. This
means that one can get from any element of the orbit to any other
element and then back to the original one by applying $\sigma$ enough
times. In the example,
\[
  \mathbf{1} \mapsto \mathbf{8} \mapsto \mathbf{6} \mapsto \mathbf{3}
  \mapsto \mathbf{2} \mapsto \mathbf{1}, \quad \mathbf{4} \mapsto
  \mathbf{7} \mapsto \mathbf{4}, \quad \mathbf{5} \mapsto \mathbf{5}.
\]
\begin{defn}{}{defn:groups2:cycle}
  A (nontrivial) \emph{cycle} is an element of $S_n$ with exactly one
  nontrivial orbit. For distinct $a_1, \dots, a_r$ in $\{1,\dots,n\}$,
  the notation $(a_1 a_2 \dots a_r)$ denotes the cycle in $S_n$ with
  nontrivial orbit $\{a_1,\dots,a_r\}$, acting as
  \[
    a_1 \mapsto a_2 \mapsto \dots \mapsto a_r \mapsto a_1.
  \]
  In this case, $r$ is the \emph{length} of the cycle. A cycle of
  length $r$ is called an \emph{$r$-cycle}.
\end{defn}
The identity is considered a cycle of length 1 in a trivial way and is
denoted by $(1)$ (and could just as well be denoted by $(i)$ for any
$i$).

Note that $(a_1 a_2 \dots a_r) = (a_2 \dots a_r a_1)$ according to the
notation introduced in \cref{defn:groups2:cycle}: the notation
determines the cycle, but a nontrivial cycle only determines the
notation `up to a cyclic permutation'.  Two cycles are \emph{disjoint}
if their nontrivial orbits are. The following observation deserves to
be highlighted, but it does not seem to deserve a proof:

\begin{lem}{}{lem:groups2:disjoint-cycles-commute}
  Disjoint cycles commute.
\end{lem}

The next one gives us the alternative notation we were looking for.

\begin{lem}{}{lem:groups2:cycle-decomposition}
  Every $\sigma \in S_n$, $\sigma \neq e$, can be written as a product
  of disjoint nontrivial cycles, in a unique way up to permutations of
  the factors.
\end{lem}

\begin{proof}
  As we have seen, every $\sigma \in S_n$ determines a partition of
  $\{1,\dots,n\}$ into orbits under the action of
  $\langle\sigma\rangle$. If $\sigma \neq e$, then
  $\langle\sigma\rangle$ has nontrivial orbits. As $\sigma$ acts as a
  cycle on each orbit, it follows that $\sigma$ may be written as a
  product of cycles.

  The proof of the uniqueness is left to the reader
  (\Cref{exer:groups2:disjoint-cycle-factor-unique}).
\end{proof}
The \emph{cycle notation} for $\sigma \in S_n$ is the (essentially)
unique expression of $\sigma$ as a product of disjoint cycles found in
\cref{lem:groups2:cycle-decomposition} (or $(1)$ for $\sigma = e$). In
our running example,
\[
  \sigma = (18632)(47),
\]
and keep in mind that this expression is unique \emph{cum grano
  salis}: we could write
\[
  \sigma = (63218)(47) = (74)(21863) = (32186)(74) = \cdots,
\]
and all of these are `the same' cycle notation for $\sigma$.

\subsection{Type and conjugacy classes in \texorpdfstring{$S_n$}{Sn}}
\label{subsec:groups2:type-conjugacy-sn}
The cycle notation has obviously annoying features---such as the
not-too-unique uniqueness pointed out a moment ago, or the fact that
$(123)$ can be an element of $S_3$ just as well as of $S_{100000}$,
and only the context can tell. However, it is invaluable as it gives
easy access to quite a bit of important information on a given
permutation. In fact, much of this information is carried already by
something even simpler than the cycle decomposition.  A
\emph{partition} of an integer $n > 0$ is a nonincreasing\footnote{Of
  course this choice is arbitrary, and nondecreasing sequences would
  do just as well.} sequence of positive integers whose sum is $n$. It
is easy to enumerate partitions for small values of $n$. For example,
$5$ has $7$ distinct partitions:
\begin{align*}
  5 &= 1 + 1 + 1 + 1 + 1\\
    &= 2 + 1 + 1 + 1\\
    &= 2 + 2 + 1\\
    &= 3 + 1 + 1\\
    &= 3 + 2\\
    &= 4 + 1\\
    &= 5.
\end{align*}
The partition $\lambda_1 \geq \lambda_2 \geq \dots \geq \lambda_r$ may
be denoted $[\lambda_1, \dots, \lambda_r]$; for example, the fourth
partition listed above would be denoted $[3, 1, 1]$. A nicer `visual'
representation is by means of the corresponding \emph{Young} (or
\emph{Ferrers}) \emph{diagram}, obtained by stacking $\lambda_1$ boxes
on top of $\lambda_2$ boxes on top of $\lambda_3$ boxes on top of
\dots. For example, the diagrams corresponding to the seven partitions
listed above are
\[ [1,1,1,1,1] \quad [2,1,1,1] \quad [2,2,1] \quad [3,1,1] \quad [3,2]
  \quad [4,1] \quad [5]
\]
(see the book for the visual Young diagrams).
\begin{defn}{}{defn:groups2:type}
  The \emph{type} of $\sigma \in S_n$ is the partition of $n$ given by
  the sizes of the orbits of the action of $\langle\sigma\rangle$ on
  $\{1,\dots,n\}$.
\end{defn}

It is hopefully clear (from the argument proving
\cref{lem:groups2:cycle-decomposition}) that the type of
$\sigma \in S_n$ is simply given by the lengths of the cycles in the
decomposition of $\sigma$ as the product of disjoint cycles, together
with as many $1$'s as needed. In our running example,
\[
  \sigma = (18632)(47) \in S_8
\]
has type $[5, 2, 1]$.  The main reason why `types' are introduced is a
consequence of the following simple observation.

\begin{lem}{}{lem:groups2:conjugate-cycle}
  Let $\tau \in S_n$, and let $(a_1 \dots a_r)$ be a cycle. Then
  \[
    \tau(a_1 \dots a_r)\tau^{-1} = (a_1\tau^{-1} \dots a_r\tau^{-1}).
  \]
\end{lem}

The funny notation $a_1\tau^{-1}$ stands for the action of the
permutation $\tau^{-1}$ on $a_1 \in \{1,\dots,n\}$; recall that we
agreed in~\Cref{subsec:groups:symmetric-groups} that we would let
our permutations act on the right, for consistency with the usual
notation for products in groups.
\begin{proof}
  This is verified by checking that both sides act in the same way on
  $\{1,\dots,n\}$. For example, for $1 \leq i < r$:
  \[
    (a_i\tau^{-1})(\tau(a_1 \dots a_r)\tau^{-1}) = a_i(a_1 \dots
    a_r)\tau^{-1} = a_{i+1}\tau^{-1}
  \]
  as it should; the other cases are left to the reader.
\end{proof}
By the usual trick of judiciously inserting identity factors
$\tau^{-1}\tau$, this formula for computing conjugates extends
immediately to any product of cycles:
\[
  \tau(a_1 \dots a_r)\cdots(b_1 \dots b_s)\tau^{-1} = (a_1\tau^{-1}
  \dots a_r\tau^{-1})\cdots(b_1\tau^{-1} \dots b_s\tau^{-1}).
\]
This holds whether the cycles are disjoint or not. However, since
$\tau$ is a bijection, disjoint cycles remain disjoint after
conjugation. This is essentially all there is to the following
important observation:
\begin{prop}{}{prop:groups2:conjugacy-type}
  Two elements of $S_n$ are conjugate in $S_n$ if and only if they
  have the same type.
\end{prop}

\begin{proof}
  The `only if' part of this statement follows immediately from the
  preceding considerations: conjugating a permutation yields a
  permutation of the same type. As for the `if' part, suppose
  \[
    \sigma_1 = (a_1 \dots a_r)(b_1 \dots b_s)\cdots(c_1 \dots c_t)
  \]
  and
  \[
    \sigma_2 = (a'_1 \dots a'_r)(b'_1 \dots b'_s)\cdots(c'_1 \dots
    c'_t)
  \]
  are two permutations with the same type, written in cycle notation,
  with $r \geq s \geq \dots \geq t$ (so the type is
  $[r, s, \dots, t]$). Let $\tau$ be any permutation such that
  $a_i = a'_i\tau$, $b_j = b'_j\tau$, \dots, $c_k = c'_k\tau$ for all
  $i$, $j$, \dots, $k$. Then \cref{lem:groups2:conjugate-cycle}
  implies $\sigma_2 = \tau\sigma_1\tau^{-1}$, so $\sigma_1$ and
  $\sigma_2$ are conjugate, as needed.
\end{proof}
\begin{exam}{}{exam:groups2:conjugate-permutations-in-s8}
  In $S_8$, $(18632)(47)$ and $(12345)(67)$ must be conjugate, since
  they have the same type. The proof of
  \cref{prop:groups2:conjugacy-type} tells us that
  $\tau(18632)(47)\tau^{-1} = (12345)(67)$ for
  \[
    \tau = \begin{pmatrix} 1 & 2 & 3 & 4 & 5 & 6 & 7 & 8 \\
      1 & 8 & 6 & 3 & 2 & 4 & 7 & 5 \end{pmatrix},
  \]
  and of course this may be checked by hand in a second. Running this
  check, and especially staring at the second row in $\tau$
  vis-\`{a}-vis the cycle notation of the first permutation, should
  clarify everything.
\end{exam}
Summarizing, the type (or the corresponding Young diagram) tells us
everything about conjugation in $S_n$.

\begin{coro}{}{coro:groups2:conjugacy-classes-count}
  The number of conjugacy classes in $S_n$ equals the number of
  partitions of $n$.
\end{coro}
For example, there are $7$ conjugacy classes in $S_5$, indexed by the
Tetris look-alikes drawn above.

It is also reasonably straightforward to compute the number of
elements in each conjugacy class, in terms of the type. For example,
in order to count the number of permutations of type $[2,2,1]$ in
$S_5$, note that there are $5! = 120$ ways to fill the corresponding
Young diagram with the numbers $1,\dots,5$:
\[
  \begin{array}{|c|c|}
    \hline
    a_1 & a_2 \\ \hline
    a_3 & a_4 \\ \hline
  \end{array}
  \quad a_5
\]
that is, $120$ ways to write a permutation as a product of two
$2$-cycles: $(a_1 a_2)(a_3 a_4)$; but switching
$a_1 \leftrightarrow a_2$ and $a_3 \leftrightarrow a_4$, as well as
switching the two cycles, gives the same permutation. Therefore there
are
\[
  \frac{120}{2 \cdot 2 \cdot 2} = 15
\]
permutations of type $[2,2,1]$. Performing this computation for all
Young diagrams for $S_5$ gives us the size of each conjugacy class,
that is, the class formula
(cf.~\Cref{subsec:groups2:center-centralizer-conjugacy-classes})
for $S_5$:
\[
  120 = 1 + 10 + 15 + 20 + 20 + 30 + 24.
\]
\begin{exam}{}{exam:groups2:no-normal-subgroup-of-order-30-in-s5}
  There are no normal subgroups of size $30$ in $S_5$. Indeed, normal
  subgroups are unions of conjugacy classes
  (\Cref{subsec:groups2:the-class-formula}); since the identity is
  in every subgroup and $30 - 1 = 29$ cannot be written as a sum of
  the numbers appearing in the class formula for $S_5$, there is no
  such subgroup.
\end{exam}

This observation will be dwarfed by much stronger results that we will
prove soon (such as \Cref{thm:groups2:alternating-group-simple-for-n-geq-5},
\Cref{coro:groups2:symmetric-group-not-solvable-for-n-geq-5}); but it is remarkable that such
precise statements can be established with so little work.

\subsection{Transpositions, parity, and the alternating group}
\label{subsec:groups2:transpositions-parity-and-the-alternating-group}
For $n \geq 1$, consider the polynomials
\[
  \Delta_n = \prod_{1 \leq i < j \leq n} (x_i - x_j) \in \Z[x_1,
  \dots, x_n],
\]
that is,
\begin{align*}
  \Delta_1 &= 1,\\
  \Delta_2 &= x_1 - x_2,\\
  \Delta_3 &= (x_1 - x_2)(x_1 - x_3)(x_2 - x_3),\\
  \Delta_4 &= (x_1 - x_2)(x_1 - x_3)(x_1 - x_4)(x_2 - x_3)
             (x_2 - x_4)(x_3 - x_4),\\
           &\phantom{{}={}}\cdots
\end{align*}
We can act with any $\sigma \in S_n$ on $\Delta_n$, by permuting the
indices according to $\sigma$:
\[
  \Delta_n\sigma := \prod_{1 \leq i < j \leq n} (x_{i\sigma} -
  x_{j\sigma}).
\]
For example,
\[
  \Delta_4(1234) = (x_2 - x_3)(x_2 - x_4)(x_2 - x_1)(x_3 - x_4) (x_3 -
  x_1)(x_4 - x_1) = -\Delta_4.
\]
In general, it is clear that $\Delta_n\sigma$ is still the product of
all binomials $(x_i - x_j)$, where the factors are permuted and some
factors may change sign in the process. Hence,
$\Delta_n\sigma = \pm\Delta_n$, where the factor $\pm 1$ depends on
$\sigma$.
\begin{defn}{}{defn:groups2:sign}
  The \emph{sign} of a permutation $\sigma \in S_n$, denoted
  $(-1)^\sigma$, is determined by the action of $\sigma$ on
  $\Delta_n$:
  \[
    \Delta_n\sigma = (-1)^\sigma \Delta_n.
  \]
  We say that a permutation is \emph{even} if its sign is $+1$ and
  \emph{odd} if its sign is $-1$.
\end{defn}

Note that $\forall\sigma,\tau \in S_n$ we have
$\Delta_n(\sigma\tau) = (\Delta_n\sigma)\tau$: it follows that
$(-1)^{\sigma\tau} = (-1)^\sigma (-1)^\tau$. Viewing $\{-1,+1\}$ as a
group under multiplication\footnote{This is the group of units in
  $\Z$; of course it is isomorphic to $C_2$.}, we see that the `sign'
function
\[
  \epsilon : S_n \to \{-1,+1\}, \quad \epsilon(\sigma) := (-1)^\sigma
\]
is a homomorphism.  Here is a different viewpoint on this sign
function. A \emph{transposition} is a cycle of length $2$. Every
permutation is a product of transpositions:

\begin{lem}{}{lem:groups2:transpositions-generate}
  Transpositions generate $S_n$.
\end{lem}

\begin{proof}
  Indeed, by \cref{lem:groups2:cycle-decomposition} it suffices to
  show that every cycle is a product of transpositions, and indeed
  \[
    (a_1 \dots a_r) = (a_1 a_2)(a_1 a_3)\cdots(a_1 a_r),
  \]
  as may be checked by applying\footnote{Don't forget that
    permutations act on the right.} both sides to every element of
  $\{1,\dots,n\}$.
\end{proof}
Of course a given permutation may be written as a product of
transpositions in many different ways. However, whether an even number
of transpositions or an odd number is needed only depends on
$\sigma$. Indeed,

\begin{lem}{}{lem:groups2:parity-well-defined}
  Let $\sigma = \tau_1 \cdots \tau_r$ be a product of
  transpositions. Then $\sigma$ is even, resp., odd, according to
  whether $r$ is even, resp., odd.
\end{lem}

\begin{proof}
  This follows immediately from the facts that $\epsilon$ is a
  homomorphism and the sign of a transposition is $-1$: indeed, $(ij)$
  acts on $\Delta_n$ by permuting its factors and changing the sign of
  an odd number of factors (for $i < j$, the factor $(x_i - x_j)$ and
  the pairs of factors $(x_i - x_k)$, $(x_k - x_j)$ for all
  $i < k < j$).
\end{proof}
\begin{defn}{}{defn:groups2:alternating-group}
  The \emph{alternating group} on $\{1,\dots,n\}$, denoted $A_n$,
  consists of all even permutations $\sigma \in S_n$.
\end{defn}

The alternating group is a normal subgroup of $S_n$, and
$[S_n : A_n] = 2$ for $n \geq 2$: indeed, $A_n = \ker\epsilon$, and
$\epsilon$ is surjective for $n \geq 2$.  It is very easy to tell
whether a permutation $\sigma$ belongs to $A_n$ in terms of the type
of $\sigma$. Indeed (by the argument proving
\cref{lem:groups2:transpositions-generate}) a cycle is even, resp.,
odd, if it has odd, resp., even, length\footnote{Note the unfortunate
  terminology clash. For example, $3$-cycles are even
  permutations.}. Pictorially, a permutation $\sigma \in S_n$ is even
if and only if $n$ and the number of rows in the Young diagram of
$\sigma$ have the same parity.

Take $n = 5$ for example: the even permutations have types
$[1,1,1,1,1]$, $[2,2,1]$, $[3,1,1]$, $[5]$ (starred types in the
book). Adding up the sizes of the corresponding conjugacy classes,
\[
  1 + 15 + 20 + 24 = 60,
\]
confirms that $A_5$ has index $2$ in $S_5$.

However, do not read too much into this computation: I am not claiming
that these are the sizes of the conjugacy classes in $A_5$, only that
these are the conjugacy classes in $S_5$ making up the normal subgroup
$A_5$. Conjugacy in $A_n$ is rather interesting and ties in nicely
with the issue of the simplicity of $A_n$.

\subsection{Conjugacy in \texorpdfstring{$A_n$}{An}; simplicity of
  \texorpdfstring{$A_n$}{An} and solvability of
  \texorpdfstring{$S_n$}{Sn}}
\label{subsec:groups2:conjugacy-an-simplicity}
Denote by $[\sigma]_{S_n}$, resp., $[\sigma]_{A_n}$, the conjugacy
class of an even permutation $\sigma$ in $S_n$, resp., $A_n$. Clearly
$[\sigma]_{A_n} \subseteq [\sigma]_{S_n}$; we proceed to compare these
two sets.
\begin{lem}{}{lem:groups2:conj-class-an-split}
  Let $n \geq 2$, and let $\sigma \in A_n$. Then
  $[\sigma]_{A_n} = [\sigma]_{S_n}$ or the size of $[\sigma]_{A_n}$ is
  half the size of $[\sigma]_{S_n}$, according to whether the
  centralizer $Z_{S_n}(\sigma)$ is not or is contained in $A_n$.
\end{lem}

\begin{proof}
  (Cf.\ \Cref{exer:groups2:idx-2-subgrp-conj-class-half}.)
  Note that $Z_{A_n}(\sigma) = A_n \cap Z_{S_n}(\sigma)$: this follows
  immediately from the definition of centralizer
  (\cref{defn:groups2:centralizer}). Now recall that the centralizer
  of $\sigma$ is its stabilizer under conjugation, and therefore the
  size of the conjugacy class of $\sigma$ equals the index of its
  centralizer.

  If $Z_{S_n}(\sigma) \subseteq A_n$, then
  $Z_{A_n}(\sigma) = Z_{S_n}(\sigma)$, so that
  \[ [S_n : Z_{S_n}(\sigma)] = [S_n : Z_{A_n}(\sigma)] = [S_n :
    A_n][A_n : Z_{A_n}(\sigma)] = 2 \cdot [A_n : Z_{A_n}(\sigma)];
  \]
  therefore, $[\sigma]_{A_n}$ is half the size of $[\sigma]_{S_n}$ in
  this case.

  If $Z_{S_n}(\sigma) \not\subseteq A_n$, then note that
  $A_n Z_{S_n}(\sigma) = S_n$: indeed, $A_n Z_{S_n}(\sigma)$ is a
  subgroup of $S_n$ (because $A_n$ is normal; cf.\
  \Cref{prop:groups:second-isomorphism-theorem}), and it properly
  contains $A_n$, so it must equal $S_n$ as $A_n$ has index $2$ in
  $S_n$. By index considerations (cf.\
  \Cref{exer:groups:hk-cardinality-formula})
  \[ [A_n : Z_{A_n}(\sigma)] = [A_n : A_n \cap Z_{S_n}(\sigma)] = [A_n
    Z_{S_n}(\sigma) : Z_{S_n}(\sigma)] = [S_n : Z_{S_n}(\sigma)],
  \]
  so the classes have the same size. Since
  $[\sigma]_{A_n} \subseteq [\sigma]_{S_n}$ in any case, it follows
  that $[\sigma]_{A_n} = [\sigma]_{S_n}$, completing the proof.
\end{proof}
Therefore, conjugacy classes of even permutations either are preserved
from $S_n$ to $A_n$ or they split into two distinct, equal-sized
classes. We are now in a position to give precise conditions
determining which happens when.

\begin{prop}{}{prop:groups2:conj-class-split-condition}
  Let $\sigma \in A_n$, $n \geq 2$. Then the conjugacy class of
  $\sigma$ in $S_n$ splits into two conjugacy classes in $A_n$
  precisely if the type of $\sigma$ consists of distinct odd numbers.
\end{prop}
\begin{proof}
  By \cref{lem:groups2:conj-class-an-split}, we have to verify that
  $Z_{S_n}(\sigma)$ is contained in $A_n$ precisely when the stated
  condition is satisfied; that is, we have to show that
  \[
    \sigma = \tau\sigma\tau^{-1} \implies \tau \text{ is even}
  \]
  precisely when the type of $\sigma$ consists of distinct odd
  numbers.

  Write $\sigma$ in cycle notation (including cycles of length $1$):
  \[
    \sigma = (a_1 \dots a_\lambda)(b_1 \dots b_\mu)\cdots(c_1 \dots
    c_\nu),
  \]
  and recall (\cref{lem:groups2:conjugate-cycle}) that
  \[
    \tau\sigma\tau^{-1} = (a_1\tau^{-1}\dots a_\lambda\tau^{-1})
    (b_1\tau^{-1}\dots b_\mu\tau^{-1})\cdots(c_1\tau^{-1}\dots
    c_\nu\tau^{-1}).
  \]
  Assume that $\lambda, \mu, \dots, \nu$ are odd and distinct. If
  $\tau\sigma\tau^{-1} = \sigma$, then conjugation by $\tau$ must
  preserve each cycle in $\sigma$, as all cycle lengths are distinct:
  \[
    \tau(a_1 \dots a_\lambda)\tau^{-1} = (a_1 \dots a_\lambda), \quad
    \text{etc.,}
  \]
  that is,
  $(a_1\tau^{-1}\dots a_\lambda\tau^{-1}) = (a_1 \dots a_\lambda)$,
  etc. This means that $\tau$ acts as a cyclic permutation on (e.g.)
  $a_1,\dots,a_\lambda$ and therefore in the same way as a power of
  $(a_1 \dots a_\lambda)$. It follows that
  \[
    \tau = (a_1 \dots a_\lambda)^r(b_1 \dots b_\mu)^s\cdots(c_1 \dots
    c_\nu)^t
  \]
  for suitable $r, s, \dots, t$. Since all cycles have odd lengths,
  each cycle is an even permutation; and $\tau$ must then be even as
  it is a product of even permutations. This proves that
  $Z_{S_n}(\sigma) \subseteq A_n$ if the stated condition holds.
  Conversely, assume that the stated condition does not hold: that is,
  either some of the cycles in the cycle decomposition have even
  length or all have odd length but two of the cycles have the same
  length.

  In the first case, let $\tau$ be an even-length cycle in the cycle
  decomposition of $\sigma$. Note that $\tau\sigma\tau^{-1} = \sigma$:
  indeed, $\tau$ commutes with itself and with all cycles in $\sigma$
  other than $\tau$.  Since $\tau$ has even length, then it is odd as
  a permutation: this shows that $Z_{S_n}(\sigma) \not\subseteq A_n$,
  as needed.

  In the second case, without loss of generality assume
  $\lambda = \mu$, and consider the odd permutation
  \[
    \tau = (a_1 b_1)(a_2 b_2)\cdots(a_\lambda b_\lambda):
  \]
  conjugating by $\tau$ simply interchanges the first two cycles in
  $\sigma$; hence $\tau\sigma\tau^{-1} = \sigma$. As $\tau$ is odd,
  this again shows that $Z_{S_n}(\sigma) \not\subseteq A_n$, and we
  are done.
\end{proof}
\begin{exam}{}{exam:groups2:a5-conjugacy-classes-from-s5-types}
  Looking again at $A_5$, we have noted
  in~\Cref{subsec:groups2:transpositions-parity-and-the-alternating-group}
  that the types of the even permutations in $S_5$ are $[1,1,1,1,1]$,
  $[2,2,1]$, $[3,1,1]$, and $[5]$.  By
  \cref{prop:groups2:conj-class-split-condition} the conjugacy classes
  corresponding to the first three types are preserved in $A_5$, while
  the last one splits.

  Therefore there are exactly $5$ conjugacy classes in $A_5$, and the
  class formula for $A_5$ is
  \[
    60 = 1 + 15 + 20 + 12 + 12.
  \]
\end{exam}
Finally! We can now complete a circle of thought begun in the first
section of this chapter. Along the way, the reader has hopefully
checked that every simple group of order $< 60$ is commutative
(\Cref{exer:groups2:no-noncomm-simple-below-60}); and we now
see why $60$ is special:

\begin{coro}{}{coro:groups2:a5-is-simple}
  The alternating group $A_5$ is a simple noncommutative group of
  order $60$.
\end{coro}

\begin{proof}
  A normal subgroup of $A_5$ is necessarily the union of conjugacy
  classes, contains the identity, and has order equal to a divisor of
  $60$ (by Lagrange's theorem). The divisors of $60$ other than $1$
  and $60$ are
  \[
    2, 3, 4, 5, 6, 10, 12, 15, 20, 30;
  \]
  counting the elements other than the identity would give one of
  \[
    1, 2, 3, 4, 5, 9, 11, 14, 19, 29
  \]
  as a sum of numbers $\neq 1$ from the class formula for $A_5$. But
  this simply does not happen.
\end{proof}
The reader will check that $A_6$ is simple, by the same method
(\Cref{exer:groups2:a6-simple-via-class-formula}). It is in
fact the case that all groups $A_n$, $n \geq 5$, are simple (and
noncommutative), implying that $S_n$ is not solvable for $n \geq 5$,
and these facts are rather important in view of applications to Galois
theory (cf., e.g., \Cref{coro:fields:solvability-by-radicals-criterion}). Note that $A_2$ is
trivial, $A_3 \cong \Z/3\Z$ is simple and abelian, and $A_4$ is not
simple (\Cref{exer:groups2:no-noncomm-simple-below-60}).

The alternating group $A_5$ is also called the \emph{icosahedral
  (rotation) group}: indeed, it is the group of symmetries of an
icosahedron obtained through rigid motions. (Can the reader
verify\footnote{It is good practice to start with smaller examples:
  for instance, the tetrahedral rotation group is isomorphic to
  $A_4$.} this fact?)  The simplicity of $A_n$ for $n \geq 5$ may be
established by studying $3$-cycles. First of all, it is natural to
wonder whether every even permutation may be written as a product of
$3$-cycles, and this is indeed so:

\begin{lem}{}{lem:groups2:an-3-cycles}
  The alternating group $A_n$ is generated by $3$-cycles.
\end{lem}

\begin{proof}
  Since every even permutation is a product of an even number of
  $2$-cycles, it suffices to show that every product of two $2$-cycles
  may be written as a product of $3$-cycles. Therefore, consider a
  product $(ab)(cd)$ with $a \neq b$, $c \neq d$. If $(ab) = (cd)$,
  then this product is the identity, and there is nothing to prove. If
  $\{a,b\}$, $\{c,d\}$ have exactly one element in common, then we may
  assume $c = a$ and observe
  \[
    (ab)(ad) = (abd).
  \]
  If $\{a,b\}$, $\{c,d\}$ are disjoint, then
  \[
    (ab)(cd) = (abc)(adc),
  \]
  and we are done.
\end{proof}
Now we can capitalize on our study of conjugacy in $A_n$:

\begin{fact}{}{fact:groups2:3-cycles-conjugacy-class}
  Let $n \geq 5$. If a normal subgroup of $A_n$ contains a $3$-cycle,
  then it contains all $3$-cycles.
\end{fact}

\begin{proof}
  Normal subgroups are unions of conjugacy classes, so we just need to
  verify that $3$-cycles form a conjugacy class in $A_n$, for
  $n \geq 5$. But they do in $S_n$, and the type of a $3$-cycle is
  $[3,1,1,\dots]$ for $n \geq 5$; hence the conjugacy class does not
  split in $A_n$, by \cref{prop:groups2:conj-class-split-condition}.
\end{proof}
The general statement now follows easily by tying up loose ends:

\begin{thm}{}{thm:groups2:alternating-group-simple-for-n-geq-5}
  The alternating group $A_n$ is simple for $n \geq 5$.
\end{thm}

\begin{proof}
  We have already checked this for $n = 5$, and the reader has checked
  it for $n = 6$. For $n > 6$, let $N$ be a nontrivial normal subgroup
  of $A_n$; we will show that necessarily $N = A_n$, by proving that
  $N$ contains $3$-cycles.

  Let $\tau \in N$, $\tau \neq (1)$, and let $\sigma \in A_n$ be a
  $3$-cycle.  Since the center of $A_n$ is trivial
  (\Cref{exer:groups2:ctr-of-an-trivial}) and $3$-cycles
  generate $A_n$, we may assume that $\tau$ and $\sigma$ do not
  commute, that is, the commutator
  \[ [\tau, \sigma] = \tau(\sigma\tau^{-1}\sigma^{-1}) =
    (\tau\sigma\tau^{-1})\sigma^{-1}
  \]
  is not the identity. This element is in $N$ (as is evident from the
  first expression, as $N$ is normal) and is a product of two
  $3$-cycles (as is evident from the second expression, since the
  conjugate of a $3$-cycle is a $3$-cycle).  Therefore, replacing
  $\tau$ by $[\tau, \sigma]$ if necessary, we may assume that
  $\tau \in N$ is a nonidentity permutation acting on $\leq 6$
  elements: that is, on a subset of a set $T \subseteq \{1,\dots,n\}$
  with $|T| = 6$.  Now we may view $A_6$ as a subgroup of $A_n$, by
  letting it act on $T$. The subgroup $N \cap A_6$ of $A_6$ is then
  normal (because $N$ is normal) and nontrivial (because
  $\tau \in N \cap A_6$ and $\tau \neq (1)$). Since $A_6$ is simple
  (\Cref{exer:groups2:a6-simple-via-class-formula}), this
  implies $N \cap A_6 = A_6$.  In particular, $N$ contains $3$-cycles.

  By \cref{fact:groups2:3-cycles-conjugacy-class}, this implies that
  $N$ contains all $3$-cycles. By \cref{lem:groups2:an-3-cycles}, it
  follows that $N = A_n$, as needed.
\end{proof}
\begin{coro}{}{coro:groups2:symmetric-group-not-solvable-for-n-geq-5}
  For $n \geq 5$, the group $S_n$ is not solvable.
\end{coro}

\begin{proof}
  Since $A_n$ is simple, the sequence
  \[
    S_n \supsetneq A_n \supsetneq \{(1)\}
  \]
  is a composition series for $S_n$. It follows that the composition
  factors of $S_n$ are $\Z/2\Z$ and $A_n$. By
  \cref{prop:groups2:solvable-equiv}, $S_n$ is not solvable.
\end{proof}

In particular, $S_5$ is a nonsolvable group of order $120$. This is in
fact the smallest order of a nonsimple, nonsolvable group; cf.\
\Cref{exer:groups2:grp-ord-lt-120-neq-60-solvable}.

\subsection*{Exercises}

\begin{exer}{}{exer:groups2:conj-class-size-example-s8}
  \exerused{} Compute the number of elements in the conjugacy class of
  \[
    \begin{pmatrix}
      1 & 2 & 3 & 4 & 5 & 6 & 7 & 8 \\
      8 & 1 & 2 & 7 & 5 & 3 & 4 & 6
    \end{pmatrix}
  \]
  in $S_8$. [\Cref{subsec:groups2:cycle-notation}]
\end{exer}
\begin{exer}{}{exer:groups2:disjoint-cycle-factor-unique}
  \exerused{} Suppose
  \[
    (a_1 \dots a_r)(b_1 \dots b_s) \cdots (c_1 \dots c_t) = (d_1 \dots
    d_u)(e_1 \dots e_v) \cdots (f_1 \dots f_w)
  \]
  are two products of disjoint cycles. Prove that the factors agree up
  to order. (Hint: The two corresponding partitions of
  $\{1, \dots, n\}$ must agree.) [\Cref{subsec:groups2:cycle-notation}]
\end{exer}
\begin{exer}{}{exer:groups2:ord-of-perm-coprime-cycle-lens}
  Assume $\sigma$ has type $[\lambda_1, \dots, \lambda_r]$ and that
  the $\lambda_i$'s are pairwise relatively prime. What is $|\sigma|$?
  What can you say about $|\sigma|$, without the additional hypothesis
  on the numbers $\lambda_i$?
\end{exer}
\begin{exer}{}{exer:groups2:partition-gen-function}
  Make sense of the `Taylor series' of the infinite product
  \[
    \frac{1}{(1-x)} \cdot \frac{1}{(1-x^2)} \cdot \frac{1}{(1-x^3)}
    \cdot \frac{1}{(1-x^4)} \cdot \frac{1}{(1-x^5)} \cdots.
  \]
  Prove that the coefficient of $x^n$ in this series is the number of
  partitions of $n$.
\end{exer}
\begin{exer}{}{exer:groups2:class-formula-sn-leq-6}
  Find the class formula for $S_n$, $n \leq 6$.
\end{exer}
\begin{exer}{}{exer:groups2:normal-subgrp-s4-ord-options}
  Let $N$ be a normal subgroup of $S_4$. Prove that $|N| = 1, 4, 12$,
  or $24$.
\end{exer}
\begin{exer}{}{exer:groups2:sn-gen-by-transposition-and-ncycle}
  \exerused{} Prove that $S_n$ is generated by $(12)$ and
  $(12 \dots n)$.  (Hint: It is enough to get all transpositions. What
  is the conjugate of $(12)$ by $(12 \dots n)$?)
  [\cref{exer:groups2:d8-in-s4-and-dihedral-in-sn},
  \Cref{subsec:fields:galois-groups-of-polynomials}]
\end{exer}
\begin{exer}{}{exer:groups2:pt-stab-in-sn-iso-s-n-minus-1}
  $\lnot$ For $n > 1$, prove that the subgroup $H$ of $S_n$ consisting
  of permutations fixing $1$ is isomorphic to $S_{n-1}$. Prove that
  there are no proper subgroups of $S_n$ properly containing $H$.
  [\ref{exer:fields:galois-grp-sn-stabilizer-no-intermediate-field}]
\end{exer}
\begin{exer}{}{exer:groups2:d8-in-s4-and-dihedral-in-sn}
  By \cref{exer:groups2:sn-gen-by-transposition-and-ncycle}, $S_4$ is
  generated by $(12)$ and $(1234)$. Prove that $(13)$ and $(1234)$
  generate a copy of $D_8$ in $S_4$. Prove that every subgroup of
  $S_4$ of order $8$ is conjugate to $\langle(13),
  (1234)\rangle$. Prove there are exactly $3$ such subgroups. For all
  $n \geq 3$ prove that $S_n$ contains a copy of the dihedral group
  $D_{2n}$, and find generators for it.
\end{exer}
\begin{exer}{}{exer:groups2:num-ncycles-and-conj-class-size-formula}
  $\lnot$
  \begin{itemize}
  \item Prove that there are exactly $(n-1)!$ $n$-cycles in $S_n$.
  \item More generally, find a formula for the size of the conjugacy
    class of a permutation of given type in $S_n$.
    [\cref{exer:groups2:sylow-count-sp-wilsons-thm}]
  \end{itemize}
\end{exer}
\begin{exer}{}{exer:groups2:sylow-count-sp-wilsons-thm}
  Let $p$ be a prime integer. Compute the number of $p$-Sylow
  subgroups of $S_p$. (Use
  \cref{exer:groups2:num-ncycles-and-conj-class-size-formula}.) Use
  this result and Sylow's third theorem to prove again the `only if'
  implication in Wilson's theorem (cf.\
  \Cref{exer:groups:wilsons-theorem}).
\end{exer}
\begin{exer}{}{exer:groups2:transitive-subgrps-s3-s4}
  \exerused{} A subgroup $G$ of $S_n$ is \emph{transitive} if the
  induced action of $G$ on $\{1, \dots, n\}$ is transitive.
  \begin{itemize}
  \item Prove that if $G \subseteq S_n$ is transitive, then $|G|$ is a
    multiple of $n$.
  \item List the transitive subgroups of $S_3$.
  \item Prove that the following subgroups of $S_4$ are all
    transitive:
    \begin{itemize}
    \item[$-$] $\langle(1234)\rangle \cong C_4$ and its conjugates,
    \item[$-$]
      $\langle(12)(34), (13)(24)\rangle \cong C_2 \times C_2$,
    \item[$-$] $\langle(12)(34), (1234)\rangle \cong D_8$ and its
      conjugates,
    \item[$-$] $A_4$, and $S_4$.
    \end{itemize}
    With a bit of stamina, you can prove that these are the
    \emph{only} transitive subgroups of $S_4$.
  \end{itemize} [\Cref{subsec:fields:galois-groups-of-polynomials}]
\end{exer}
\begin{exer}{}{exer:groups2:sign-eq-det-of-perm-matrix}
  (If you know about determinants.) Prove that the sign of a
  permutation $\sigma$, as defined in \cref{defn:groups2:sign}, equals
  the determinant of the matrix $M_\sigma$ defined in
  \Cref{exer:groups:permutation-matrix}.
\end{exer}
\begin{exer}{}{exer:groups2:ctr-of-an-trivial}
  \exerused{} Prove that the center of $A_n$ is trivial for
  $n \geq 4$.  [\Cref{subsec:groups2:conjugacy-an-simplicity}]
\end{exer}
\begin{exer}{}{exer:groups2:justify-pictorial-parity-recipe}
  Justify the `pictorial' recipe given in
  \Cref{subsec:groups2:transpositions-parity-and-the-alternating-group}
  to decide whether a permutation is even.
\end{exer}
\begin{exer}{}{exer:groups2:num-conj-classes-in-an-list}
  The number of conjugacy classes in $A_n$, $n \geq 2$, is (allegedly)
  \[
    1, 3, 4, 5, 7, 9, 14, 18, 24, 31, 43, \dots
  \]
  Check the first several numbers in this list by finding the class
  formulas for the corresponding alternating groups.
\end{exer}
\begin{exer}{}{exer:groups2:class-formula-a4-no-subgrp-ord-6}
  \exerused{}
  \begin{itemize}
  \item Find the class formula for $A_4$.
  \item Use it to prove that $A_4$ has no subgroup of order $6$.
    [\Cref{subsec:groups:the-index-and-lagrange-s-theorem}]
  \end{itemize}
\end{exer}
\begin{exer}{}{exer:groups2:an-proper-subgrp-idx-geq-n}
  For $n \geq 5$, let $H$ be a proper subgroup of $A_n$. Prove that
  $[A_n : H] \geq n$. Prove that $A_n$ does have a subgroup of index
  $n$ for all $n \geq 3$.
\end{exer}
\begin{exer}{}{exer:groups2:an-no-action-on-small-set}
  Prove that for $n \geq 5$ there are no nontrivial actions of $A_n$
  on any set $S$ with $|S| < n$. Construct\footnote{You can think
    algebraically if you want; if you prefer geometry, visualize pairs
    of opposite edges on a tetrahedron.} a nontrivial action of $A_4$
  on a set $S$, $|S| = 3$. Is there a nontrivial action of $A_4$ on a
  set $S$ with $|S| = 2$?
\end{exer}
\begin{exer}{}{exer:groups2:a5-ord-2-elts-and-sylow-count}
  $\lnot$ Find all fifteen elements of order $2$ in $A_5$, and prove
  that $A_5$ has exactly five $2$-Sylow subgroups.
  [\cref{exer:groups2:a5-unique-simple-grp-ord-60}]
\end{exer}
\begin{exer}{}{exer:groups2:a6-simple-via-class-formula}
  \exerused{} Prove that $A_6$ is simple, by using its class formula
  (as is done for $A_5$ in the proof of
  \cref{coro:groups2:a5-is-simple}). [\Cref{subsec:groups2:conjugacy-an-simplicity}]
\end{exer}
\begin{exer}{}{exer:groups2:a5-unique-simple-grp-ord-60}
  $\lnot$ Verify that $A_5$ is the \emph{only} simple group of order
  $60$, up to isomorphism. (Hint: By
  \Cref{exer:groups2:simple-grp-ord-60-subgrp-idx-5}, a simple group
  $G$ of order $60$ contains a subgroup of index $5$. Use this fact to
  construct a homomorphism $G \to S_5$, and prove that the image of
  this homomorphism must be $A_5$.) Note that $A_5$ has exactly five
  $2$-Sylow subgroups; cf.\
  \cref{exer:groups2:a5-ord-2-elts-and-sylow-count}. Thus, the other
  possibility contemplated in
  \Cref{exer:groups2:simple-grp-ord-60-subgrp-idx-5} does not
  occur. [\ref{exer:groups2:simple-grp-ord-60-subgrp-idx-5}]
\end{exer}

\section{Products of groups}
\label{sec:groups2:products-of-groups}

We already know that products exist in $\categ{Grp}$
(see~\Cref{subsec:groups:products-et-al}); here we analyze this
notion further and explore variations on the same theme, with an eye
towards the question of determining the information needed to
reconstruct a group from its composition factors.

\subsection{The direct product}
\label{subsec:groups2:the-direct-product}
Recall from~\Cref{subsec:groups:products-et-al} that the
\emph{(direct) product} of two groups $H$, $K$ is the group supported
on the set $H \times K$, with operation defined componentwise. We have
checked (\Cref{prop:groups:direct-product-universal}) that the direct
product satisfies the universal property defining products in the
category $\categ{Grp}$.  There are situations in which the direct
product of two subgroups $N$, $H$ of a group $G$ may be realized as a
subgroup of $G$. Recall
(\Cref{prop:groups:second-isomorphism-theorem}) that if one of the
subgroups is normal, then the subset $NH$ of $G$ is in fact a subgroup
of $G$. The relation between $NH$ and $N \times H$ depends on how $N$
and $H$ intersect in $G$, so we take a look at this intersection.  The
`commutator' $[A, B]$ of two subsets $A$, $B$ of $G$
(see~\Cref{subsec:groups2:the-commutator-subgroup-derived-series-and-solvability})
is the subgroup generated by all commutators $[a, b]$ with $a \in A$,
$b \in B$.
\begin{lem}{}{lem:groups2:commutator-of-normal-subgroups-in-intersection}
  Let $N$, $H$ be normal subgroups of a group $G$. Then
  $[N, H] \subseteq N \cap H$.
\end{lem}
\begin{proof}
  It suffices to verify this on generators; that is, it suffices to
  check that
  \[ [n, h] = n(hn^{-1}h^{-1}) = (nhn^{-1})h^{-1} \in N \cap H
  \]
  for all $n \in N$, $h \in H$. But the first expression and the
  normality of $N$ show that $[n, h] \in N$; the second expression and
  the normality of $H$ show that $[n, h] \in H$.
\end{proof}
\begin{coro}{}{coro:groups2:disjoint-normal-subgroups-commute}
  Let $N$, $H$ be normal subgroups of a group $G$. Assume
  $N \cap H = \{e\}$. Then $N$, $H$ commute with each other:
  \[
    (\forall n \in N)\,(\forall h \in H) \quad nh = hn.
  \]
\end{coro}
\begin{proof}
  By \cref{lem:groups2:commutator-of-normal-subgroups-in-intersection}, $[N, H] = \{e\}$ if $N \cap H = \{e\}$;
  the result follows immediately.
\end{proof}
In fact, under the same hypothesis more is true:

\begin{prop}{}{prop:groups2:disjoint-normal-subgroups-product-is-direct}
  Let $N$, $H$ be normal subgroups of a group $G$, such that
  $N \cap H = \{e\}$. Then $NH \cong N \times H$.
\end{prop}
\begin{proof}
  Consider the function
  \[\varphi : N \times H \to NH\]
  defined by $\varphi(n, h) = nh$. Under the stated hypothesis,
  $\varphi$ is a group homomorphism: indeed
  \begin{align*}
    \varphi((n_1, h_1) \cdot (n_2, h_2))
    &= \varphi((n_1 n_2, h_1 h_2)) \\
    &= n_1 n_2 h_1 h_2 \\
    &= n_1 h_1 n_2 h_2
    && \text{since $N$, $H$ commute by \cref{coro:groups2:disjoint-normal-subgroups-commute}} \\
    &= \varphi((n_1, h_1)) \cdot \varphi((n_2, h_2)).
  \end{align*}
  The homomorphism $\varphi$ is surjective by definition of $NH$. To
  verify it is injective, consider its kernel:
  \[\ker\varphi = \{(n, h) \in N \times H \mid nh = e\}.\]
  If $nh = e$, then $n \in N$ and $n = h^{-1} \in H$; thus $n = e$
  since $N \cap H = \{e\}$. Using the same token for $h$, we conclude
  $h = e$; hence $(n, h)$ is the identity in $N \times H$, proving
  that $\varphi$ is injective.

  Thus $\varphi$ is an isomorphism, as needed.
\end{proof}
\begin{rem}{}{rem:groups2:order-pq-cyclic-alternative-argument}
  This result gives an alternative argument for the proof of
  \Cref{fact:groups2:order-pq-cyclic}: if $|G| = pq$, with $p < q$
  prime integers, and $G$ contains normal subgroups $H$, $K$ of order
  $p$, $q$, respectively (as is the case if $q \not\equiv 1 \bmod p$,
  by Sylow), then $H \cap K = \{e\}$ necessarily, and then
  \cref{prop:groups2:disjoint-normal-subgroups-product-is-direct} shows $HK \cong H \times K$. As
  $|HK| = |G| = pq$, this proves
  $G \cong H \times K \cong \Z/p\Z \times \Z/q\Z$. Finally, $(1, 1)$
  has order $pq$ in this group, so $G$ is cyclic, with the same
  conclusion we obtained in \Cref{fact:groups2:order-pq-cyclic}.
\end{rem}

\subsection{Exact sequences of groups; extension problem}
\label{subsec:groups2:exact-sequences-of-groups-extension-problem}
Of course, the hypothesis that both subgroups $N$, $H$ are normal is
necessary for the result of \cref{prop:groups2:disjoint-normal-subgroups-product-is-direct}: for example, the
permutations $(123)$ and $(12)$ generate subgroups $N$, $H$ of $S_3$
meeting only at $\{e\}$, and $N$ is normal in $S_3$, but $S_3 = NH$ is
\emph{not} isomorphic to the direct product of $N$ and $H$. It is
natural to examine this more general situation.  Let $N$, $H$ be
subgroups of a group $G$, with $N$ normal (but with no \emph{a priori}
assumptions on $H$) and such that $N \cap H = \{e\}$; assume $G =
NH$. We are after a description of the structure of $G$ in terms of
the structure of $N$ and $H$.  It is notationally convenient to use
the language of \emph{exact sequences}, introduced for modules
in~\Cref{subsec:rings:complexes-and-exact-sequences}. A (short)
exact sequence of groups is a sequence of groups and group
homomorphisms
\[1 \longrightarrow N \xrightarrow{\varphi} G \xrightarrow{\psi} H
  \longrightarrow 1\] where $\psi$ is surjective and $\varphi$
identifies $N$ with $\ker\psi$. In other words (by the first
isomorphism theorem), use $\varphi$ to identify $N$ with a subgroup of
$G$; then the sequence is exact if $N$ is normal in $G$ and $\psi$
induces an isomorphism $G/N \to H$.  The reader should pause a moment
and check that if $G$, $N$, $H$ are abelian, then this notion matches
precisely the notion of short exact sequence of abelian groups (i.e.,
$\Z$-modules) from
\Cref{subsec:rings:complexes-and-exact-sequences}; a notational
difference is that here the trivial group is denoted\footnote{This is
  not unreasonable, since groups are more often written
  `multiplicatively' rather than additively, so the identity element
  is more likely to be denoted $1$ rather than $0$.}  `$1$' rather
than `$0$'.

Of course there always is an exact sequence
\[1 \longrightarrow N \longrightarrow N \times H \longrightarrow H
  \longrightarrow 1\,:\] map $n \in N$ to $(n, e_H)$ and
$(n, h) \in N \times H$ to $h$.  However, keep in mind that this is a
very special case: to reiterate the example mentioned above, there
also is an exact sequence
\[1 \longrightarrow C_3 \longrightarrow S_3 \longrightarrow C_2
  \longrightarrow 1\,,\] yet $S_3 \not\cong C_3 \times C_2$.
\begin{defn}{}{defn:groups2:group-extension}
  Let $N$, $H$ be groups. A group $G$ is an \emph{extension of $H$ by
    $N$} if there is an exact sequence of groups
  \[1 \longrightarrow N \longrightarrow G \longrightarrow H
    \longrightarrow 1\,.\]
\end{defn}
The \emph{extension problem} aims to describe all extensions of two
given groups, up to isomorphism. For example, there are two extensions
of $C_2$ by $C_3$: namely $C_6 \cong C_3 \times C_2$ and $S_3$; we
will soon be able to verify that, up to isomorphism, there are no
other extensions.  The extension problem is the `second half' of the
classification problem: the first half consists of determining all
simple groups, and the second half consists of figuring out how these
can be put together to construct any group.\footnote{As mentioned
  earlier, the first half has been settled, although the complexity of
  the work leading to its solution justifies some skepticism
  concerning the absolute correctness of the proof. The status of the
  second half is, as far as I know, (even) murkier.} For example, if
\[G = G_0 \supsetneq G_1 \supsetneq G_2 \supsetneq G_3 \supsetneq G_4
  = \{e\}\] is a composition series, with (simple) quotients
$H_i = G_i/G_{i+1}$, then $G$ is an extension of $H_0$ by an extension
of $H_1$ by an extension of $H_2$ by $H_3$: knowing the composition
factors of $G$ and the extension process, it should in principle be
possible to reconstruct $G$.

We are going to `solve' the extension problem in the particular case
in which $H$ is also a subgroup of $G$, intersecting $N$ at $\{e\}$.
\begin{defn}{}{defn:groups2:split-extension}
  An exact sequence of groups
  \[1 \longrightarrow N \xrightarrow{\varphi} G \xrightarrow{\psi} H
    \longrightarrow 1\] (or the corresponding extension) is said to
  \emph{split} if $\psi$ has a right inverse.

  In this case $H$ may be identified with a subgroup of $G$ so that
  $N \cap H = \{e\}$.
\end{defn}

We encountered this terminology
in~\Cref{subsec:rings:split-exact-sequences} for modules, thus for
abelian groups. Note that the notion examined there appears to be more
restrictive than \cref{defn:groups2:split-extension}, since it requires $G$ to be
isomorphic to a direct product $N \times H$. This apparent mismatch
evaporates because of \cref{prop:groups2:disjoint-normal-subgroups-product-is-direct}: in the abelian case,
every split extension (according to \cref{defn:groups2:split-extension}) is in
fact a direct product. Of course split extensions are anyway very
special, since quotients of a group $G$ are usually not isomorphic to
subgroups of $G$, even in the abelian case (cf.\
\cref{exer:groups2:z-mult-2-seq-exact-not-split}).
\begin{lem}{}{lem:groups2:complement-gives-split-extension}
  Let $N$ be a normal subgroup of a group $G$, and let $H$ be a
  subgroup of $G$ such that $G = NH$ and $N \cap H = \{e\}$. Then $G$
  is a split extension of $H$ by $N$.
\end{lem}
\begin{proof}
  We have to construct an exact sequence
  \[1 \longrightarrow N \longrightarrow G \longrightarrow H
    \longrightarrow 1\,;\] we let $N \to G$ be the inclusion map, and
  we prove that $G/N \cong H$. For this, consider the composition
  \[\alpha : H \hookrightarrow G \twoheadrightarrow G/N.\]
  Then $\alpha$ is surjective: indeed, since $G = NH$,
  $\forall g \in G$ we have $g = nh$ for some $n \in N$ and $h \in H$,
  and then
  \[gN = nhN = h(h^{-1}nh)N = hN = \alpha(h).\] Further,
  $\ker\alpha = \{h \in H \mid hN = N\} = N \cap H = \{e\}$; therefore
  $\alpha$ is also injective, as needed.
\end{proof}
To recap, if in the situation of \cref{lem:groups2:complement-gives-split-extension} we also
require that $H$ be normal in $G$, then $G$ is necessarily isomorphic
to the `trivial' extension $N \times H$: this is what we have proved
in \cref{prop:groups2:disjoint-normal-subgroups-product-is-direct}. We are seeking to describe the extension
`even if' $H$ is not normal in $G$.

\subsection{Internal/semidirect products}
\label{subsec:groups2:internal-semidirect-products}
The attentive reader should have noticed that the key to
\cref{prop:groups2:disjoint-normal-subgroups-product-is-direct} is really \cref{coro:groups2:disjoint-normal-subgroups-commute}: if both $N$
and $H$ are normal and $N \cap H = \{e\}$, then $N$ and $H$ commute
with each other. This is what ultimately causes the extension $NH$ to
be trivial. Now, recall that as soon as $N$ is normal, then every
subgroup $H$ of $G$ acts on $N$ by conjugation: in fact (cf.\
\Cref{exer:groups2:conj-action-kernel-is-centralizer}) conjugation
determines a homomorphism
\[\gamma : H \to \Aut_{\categ{Grp}}(N), \quad
  h \mapsto \gamma_h.\] (Explicitly, for $h \in H$ the automorphism
$\gamma_h : N \to N$ acts by $\gamma_h(n) := hnh^{-1}$.) The subgroups
$H$ and $N$ commute precisely when $\gamma$ is
trivial. \Cref{coro:groups2:disjoint-normal-subgroups-commute} shows that if $N$ and $H$ are both
normal and $N \cap H = \{e\}$, then $\gamma$ is indeed trivial.  This
is the crucial remark. The next several considerations may be
summarized as follows: if $N$ is normal in $G$, $H$ is a subgroup of
$G$, $N \cap H = \{e\}$ and $G = NH$, then the extension $G$ of $H$ by
$N$ may be reconstructed from the conjugation action
$\gamma : H \to \Aut_{\categ{Grp}}(N)$. The reader is advised to stare
at the following triviality, which is the motivating observation for
the general discussion:
\begin{equation}
  \tag{*}
  (\forall n_1, n_2 \in N),\, (\forall h_1, h_2 \in H) \quad
  n_1 h_1 n_2 h_2 = (n_1(h_1 n_2 h_1^{-1}))(h_1 h_2).
\end{equation}
This says that if we know the conjugation action of $H$ on $N$, then
we can recover the operation in $G$ from this information and from the
operations in $N$ and $H$.

Here is the general discussion. It is natural to abstract the
situation and begin with any two groups $N$, $H$ and an arbitrary
homomorphism\footnote{The reason why I am not denoting the image of
  $h$ by $\theta$ as $\theta(h)$ is that this is an automorphism of
  $N$ and I dislike the notation $\theta(h)(n)$ for the image of
  $n \in N$ obtained by applying the automorphism corresponding to
  $h$. The alternative $\theta_h(n)$ looks a little easier to parse.}
\[\theta : H \to \Aut_{\categ{Grp}}(N), \quad
  h \mapsto \theta_h.\] Define an operation $\bullet_\theta$ on the
set $N \times H$ as follows: for $n_1, n_2 \in N$ and
$h_1, h_2 \in H$, let
\[(n_1, h_1) \bullet_\theta (n_2, h_2) := (n_1 \theta_{h_1}(n_2),\,
  h_1 h_2).\] This will look more reasonable once it is compared with
$(*)$!
\begin{lem}{}{lem:groups2:semidirect-product-operation-is-a-group}
  The resulting structure $(N \times H, \bullet_\theta)$ is a group,
  with identity element $(e_N, e_H)$.
\end{lem}
\begin{proof}
  The reader should carefully verify this. For example, inverses exist
  because
  \begin{align*}
    (n_1, h_1) \bullet_\theta (\theta_{h_1^{-1}}(n_1^{-1}), h_1^{-1})
    &= (n_1 \theta_{h_1}(\theta_{h_1^{-1}}(n_1^{-1})),\, h_1 h_1^{-1})
      = (n_1 n_1^{-1},\, e_H) = (e_N, e_H)
  \end{align*}
  and similarly in the reverse order.
\end{proof}
\begin{defn}{}{defn:groups2:semidirect-product}
  The group $(N \times H, \bullet_\theta)$ is a \emph{semidirect
    product} of $N$ and $H$ and is denoted by $N \rtimes_\theta H$.
\end{defn}
For example, the ordinary direct product is a semidirect product and
corresponds to $\theta = $ the trivial map. If the reader feels a
little uneasy about giving one name (semidirect product) for a whole
host of different groups supported on the Cartesian product, welcome
to the club. In fact, it gets worse still: it is not uncommon to omit
`$\theta$' from the notation and simply write $N \rtimes H$ for a
semidirect product.\footnote{This is actually OK, if $N$ and $H$ are
  given as subgroups of a common group $G$, in which case the implicit
  action is just conjugation.}

In any case, the notation $\bullet_\theta$ is too heavy to carry
around, so we generally revert back to the usual simple juxtaposition
of elements in order to denote multiplication in $N \rtimes_\theta H$.
The following proposition checks that semidirect products are split
extensions:

\begin{prop}{}{prop:groups2:semidirect-product-properties}
  Let $N$, $H$ be groups, and let
  $\theta : H \to \Aut_{\categ{Grp}}(N)$ be a homomorphism; let
  $G = N \rtimes_\theta H$ be the corresponding semidirect
  product. Then
  \begin{itemize}
  \item $G$ contains isomorphic copies of $N$ and $H$;
  \item the natural projection $G \to H$ is a surjective homomorphism,
    with kernel $N$; thus $N$ is normal in $G$, and the sequence
    \[1 \longrightarrow N \longrightarrow N \rtimes_\theta H
      \longrightarrow H \longrightarrow 1\] is (split) exact;
  \item $N \cap H = \{e_G\}$;
  \item $G = NH$;
  \item the homomorphism $\theta$ is realized by conjugation in $G$:
    that is, for $h \in H$ and $n \in N$ we have
    $\theta_h(n) = hnh^{-1}$ in $G$.
  \end{itemize}
\end{prop}
\begin{proof}
  The functions $N \to G$, $H \to G$ defined for $n \in N$, $h \in H$
  by
  \[n \mapsto (n, e_H), \qquad h \mapsto (e_N, h)\] are manifestly
  injective homomorphisms, allowing us to identify $N$, $H$ with the
  corresponding subgroups of $G$. It is clear that
  $N \cap H = \{(e_N, e_H)\} = \{e_G\}$, and
  \[(n, e_H) \bullet_\theta (e_N, h) = (n, h)\] shows that $G = NH$.

  The projection $G \to H$ defined by $(n, h) \mapsto h$ is a
  surjective homomorphism, with kernel $N$; therefore $N$ is normal in
  $G$. Finally,
  \[(e_N, h) \bullet_\theta (n, e_H) \bullet_\theta (e_N, h)^{-1} =
    (\theta_h(n), h) \bullet_\theta (e_N, h^{-1}) = (\theta_h(n),
    e_H),\] as claimed in the last point.
\end{proof}
Our original goal of `reconstructing' a given split extension of a
group $H$ by a group $N$ is a sort of converse to this proposition.
More precisely,

\begin{prop}{}{prop:groups2:internal-semidirect-product-characterization}
  Let $N$, $H$ be subgroups of a group $G$, with $N$ normal in $G$.
  Assume that $N \cap H = \{e\}$, and $G = NH$. Let
  $\gamma : H \to \Aut_{\categ{Grp}}(N)$ be defined by conjugation:
  for $h \in H$, $n \in N$,
  \[\gamma_h(n) = hnh^{-1}.\]
  Then $G \cong N \rtimes_\gamma H$.
\end{prop}
\begin{proof}
  Define a function
  \[\varphi : N \rtimes_\gamma H \to G\]
  by $\varphi(n, h) = nh$; this is clearly a bijection. We need to
  verify that $\varphi$ is a homomorphism, and indeed
  $(\forall n_1, n_2 \in N)$, $(\forall h_1, h_2 \in H)$:
  \begin{align*}
    \varphi((n_1, h_1) \bullet_\gamma (n_2, h_2))
    &= \varphi((n_1 \gamma_{h_1}(n_2),\, h_1 h_2)) \\
    &= \varphi((n_1(h_1 n_2 h_1^{-1}),\, h_1 h_2)) \\
    &= n_1 h_1 n_2 (h_1^{-1} h_1) h_2 = (n_1 h_1)(n_2 h_2) \\
    &= \varphi((n_1, h_1))\,\varphi((n_2, h_2))
  \end{align*}
  as needed.
\end{proof}
When realized `within' a group, as in the previous proposition, the
semidirect product is sometimes called an \emph{internal} product.
\begin{rem}{}{rem:groups2:trivial-action-gives-direct-product}
  If $N$ and $H$ commute, then the conjugation action of $H$ on $N$ is
  trivial; therefore $\gamma$ is the trivial map, and the semidirect
  product $N \rtimes_\gamma H$ is the direct product $N \times H$.
  Thus, \cref{prop:groups2:internal-semidirect-product-characterization} recovers the result of
  \cref{prop:groups2:disjoint-normal-subgroups-product-is-direct} in this case.
\end{rem}
\begin{exam}{}{exam:groups2:automorphism-group-of-c3}
  The automorphism group of $C_3$ is isomorphic to the cyclic group
  $C_2$: if $C_3 = \{e, y, y^2\}$, then the two automorphisms of $C_3$
  are
  \[
    \id : \begin{cases} e \mapsto e, \\ y \mapsto y, \\ y^2 \mapsto
      y^2,
    \end{cases}
    \qquad \sigma : \begin{cases} e \mapsto e, \\ y \mapsto y^2, \\
      y^2 \mapsto y.
    \end{cases}
  \]
  Therefore, there are two homomorphisms
  $C_2 \to \Aut_{\categ{Grp}}(C_3)$: the trivial map, and the
  isomorphism sending the identity to $\id$ and the nonidentity
  element to $\sigma$. The semidirect product corresponding to the
  trivial map is the direct product $C_3 \times C_2 \cong C_6$; the
  other semidirect product $C_3 \rtimes C_2$ is isomorphic to
  $S_3$. This can of course be checked by hand (and you should do it,
  for fun); but it also follows immediately from
  \cref{prop:groups2:internal-semidirect-product-characterization}, since $N = \langle(123)\rangle$,
  $H = \langle(12)\rangle \subseteq S_3$ satisfy the hypotheses of
  this result.
\end{exam}
The reader should contemplate carefully the slightly more general case
of dihedral groups (\cref{exer:groups2:d2n-as-semidirect-prod}); this
enriches (and in a sense explains) the discussion presented in
\Cref{fact:groups2:order-2q-dihedral}.

In fact, semidirect products shed light on all groups of order $pq$,
for $p < q$ primes; the reader should be able to complete the
classification of these groups begun in
\Cref{fact:groups2:order-pq-cyclic} and show that if
$q \equiv 1 \bmod p$, then there is exactly one such noncommutative
group up to isomorphism (\cref{exer:groups2:classify-grps-ord-pq}).

\subsection*{Exercises}

The reader would in fact be well-advised to try to use semidirect
products to classify groups of small order: if a nontrivial normal
subgroup $N$ is found (typically by applying Sylow's theorems), with
some luck the classification is reduced to the study of possible
homomorphisms from known groups to $\Aut_{\categ{Grp}}(N)$ and can be
carried out. See \cref{exer:groups2:classify-grps-ord-28} for a
further illustration of this technique.
\begin{exer}{}{exer:groups2:all-sylow-normal-implies-direct-product}
  \exerused{} Let $G$ be a finite group, and let $P_1, \dots, P_r$ be
  its nontrivial Sylow subgroups. Assume all $P_i$ are normal in $G$.
  \begin{itemize}
  \item Prove that $G \cong P_1 \times \dots \times P_r$. (Induction
    on $r$; use \cref{prop:groups2:disjoint-normal-subgroups-product-is-direct}.)
  \item Prove that $G$ is nilpotent. (Hint: Mod out by the center, and
    work by induction on $|G|$. What is the center of a direct product
    of groups?)
  \end{itemize}
  Together with
  \cref{exer:groups2:nilpotent-grp-upper-central-series}, this shows
  that a finite group is nilpotent if and only if each of its Sylow
  subgroups is
  normal. [\cref{exer:groups2:nilp-grp-proper-subgrp-lt-normalizer},
  \Cref{subsec:groups2:classification-of-finite-abelian-groups}]
\end{exer}
\begin{exer}{}{exer:groups2:comp-factors-of-extension}
  Let $G$ be an extension of $H$ by $N$. Prove that the composition
  factors of $G$ are the collection of the composition factors of $H$
  and those of $N$.
\end{exer}
\begin{exer}{}{exer:groups2:normal-series-as-exact-sequences}
  Let
  \[G = G_0 \supsetneq G_1 \supsetneq \dots \supsetneq G_r = \{e\}\]
  be a normal series. Show how to `connect' $\{e\}$ to $G$ by means of
  $r$ exact sequences of groups using the groups $G_i$ and the
  quotients $H_i = G_i/G_{i+1}$.
\end{exer}
\begin{exer}{}{exer:groups2:z-mult-2-seq-exact-not-split}
  \exerused{} Prove that the sequence
  \[0 \longrightarrow \Z \xrightarrow{\cdot 2} \Z \longrightarrow
    \Z/2\Z \longrightarrow 0\] is exact but does not
  split. [\Cref{subsec:groups2:exact-sequences-of-groups-extension-problem}]
\end{exer}
\begin{exer}{}{exer:groups2:split-seq-of-grps-left-inverse-q}
  In \Cref{prop:rings:split-sequence} we have seen that if an exact
  sequence
  \[0 \longrightarrow M \xrightarrow{\varphi} N \longrightarrow
    N/(\varphi(M)) \longrightarrow 0\] of abelian groups splits, then
  $\varphi$ has a left-inverse. Is this necessarily the case for split
  sequences of groups?
\end{exer}
\begin{exer}{}{exer:groups2:semidirect-prod-is-grp}
  Prove \cref{lem:groups2:semidirect-product-operation-is-a-group}.
\end{exer}
\begin{exer}{}{exer:groups2:automorphism-realized-as-conj}
  Let $N$ be a group, and let $\alpha : N \to N$ be an automorphism of
  $N$. Prove that $\alpha$ may be realized as conjugation, in the
  sense that there exists a group $G$ containing $N$ as a normal
  subgroup and such that $\alpha(n) = gng^{-1}$ for some $g \in G$.
\end{exer}
\begin{exer}{}{exer:groups2:semidirect-of-solvable-is-solvable}
  Prove that any semidirect product of two solvable groups is
  solvable.  Show that semidirect products of nilpotent groups need
  not be nilpotent.
\end{exer}
\begin{exer}{}{exer:groups2:comm-semidirect-prod-is-direct}
  \exerused{} Prove that if $G = N \rtimes H$ is commutative, then
  $G \cong N \times
  H$. [\Cref{subsec:groups2:classification-of-finite-abelian-groups}]
\end{exer}
\begin{exer}{}{exer:groups2:coprime-normal-subgrp-semidirect}
  Let $N$ be a normal subgroup of a finite group $G$, and assume that
  $|N|$ and $|G/N|$ are relatively prime. Assume there is a subgroup
  $H$ in $G$ such that $|H| = |G/N|$. Prove that $G$ is a semidirect
  product of $N$ and $H$.
\end{exer}
\begin{exer}{}{exer:groups2:d2n-as-semidirect-prod}
  \exerused{} For all $n > 0$ express $D_{2n}$ as a semidirect product
  $C_n \rtimes_\theta C_2$, finding $\theta$
  explicitly. [\Cref{subsec:groups2:internal-semidirect-products}]
\end{exer}
\begin{exer}{}{exer:groups2:classify-grps-ord-pq}
  \exerused{} Classify groups $G$ of order $pq$, with $p < q$ prime:
  show that if $|G| = pq$, then either $G$ is cyclic or
  $q \equiv 1 \bmod p$ and there is exactly one isomorphism class of
  noncommutative groups of order $pq$ in this case. (You will likely
  have to use the fact that $\Aut_{\categ{Grp}}(C_q) \cong C_{q-1}$ if
  $q$ is prime; cf.\ \Cref{exer:groups:aut-z-aut-cp}.) [\Cref{subsec:groups2:applications},
  \Cref{subsec:groups2:internal-semidirect-products}]
\end{exer}
\begin{exer}{}{exer:groups2:semidirect-ker-theta-is-action-ker}
  $\lnot$ Let $G = N \rtimes_\theta H$ be a semidirect product, and
  let $K$ be the subgroup of $G$ corresponding to
  $\ker\theta \subseteq H$.  Prove that $K$ is the kernel of the
  action of $G$ on the set $G/H$ of left-cosets of
  $H$. [\cref{exer:groups2:c2xc2-semidirect-s3-iso-s4}]
\end{exer}
\begin{exer}{}{exer:groups2:c2xc2-semidirect-s3-iso-s4}
  Recall that $S_3 \cong \Aut_{\categ{Grp}}(C_2 \times C_2)$
  (\Cref{exer:groups:aut-c2xc2-iso-s3}). Let $\iota$ be this
  isomorphism. Prove that
  $(C_2 \times C_2) \rtimes_\iota S_3 \cong S_4$. (Hint:
  \cref{exer:groups2:semidirect-ker-theta-is-action-ker}.)
\end{exer}
\begin{exer}{}{exer:groups2:classify-grps-ord-28}
  \exerused{} Let $G$ be a group of order $28$.
  \begin{itemize}
  \item Prove that $G$ contains a normal subgroup $N$ of order $7$.
  \item Recall (or prove again) that, up to isomorphism, the only
    groups of order $4$ are $C_4$ and $C_2 \times C_2$. Prove that
    there are two homomorphisms $C_4 \to \Aut_{\categ{Grp}}(N)$ and
    two homomorphisms $C_2 \times C_2 \to \Aut_{\categ{Grp}} (N)$ up
    to the choice of generators for the sources.
  \item Conclude that there are four groups of order $28$ up to
    isomorphism: the two direct products $C_4 \times C_7$,
    $C_2 \times C_2 \times C_7$, and two noncommutative groups.
  \item Prove that $D_{28} \cong C_2 \times D_{14}$. The other
    noncommutative group of order $28$ is a \emph{generalized
      quaternionic group}.
  \end{itemize}
  [\Cref{subsec:groups2:internal-semidirect-products}]
\end{exer}
\begin{exer}{}{exer:groups2:q8-not-semidirect-prod}
  Prove that the quaternionic group $Q_8$ (cf.\
  \Cref{exer:rings:quaternions}) cannot be written as a semidirect
  product of two nontrivial subgroups.
\end{exer}
\begin{exer}{}{exer:groups2:quaternions-as-semidirect-prod}
  Prove that the multiplicative group $\field{H}^*$ of nonzero
  quaternions (cf.\ \Cref{exer:rings:quaternions}) is isomorphic to a
  semidirect product $\mathrm{SU}(2) \rtimes \R^+$. (Hint:
  \Cref{exer:rings:quaternion-norm}.) Is this semidirect product in
  fact direct?
\end{exer}

\section{Finite abelian groups}
\label{sec:groups2:finite-abelian-groups}

I will end this chapter by treating in some detail the classification
theorem for finite abelian groups mentioned
in~\Cref{subsec:groups:example-subgroup-generated-by-a-subset}.

\subsection{Classification of finite abelian groups}
\label{subsec:groups2:classification-of-finite-abelian-groups}
Now that we have acquired more familiarity with products, we are in a
position to classify all finite \emph{abelian} groups.\footnote{Of
  course fancier semidirect products will not be needed here; cf.\
  \cref{exer:groups2:comm-semidirect-prod-is-direct}.}

In due time (\Cref{prop:linalg:euclidean-domain-normal-form},
\Cref{exer:linalg:fin-gen-abel-grp-direct-sum-cyclic-via-matrices}) we
will in fact be able to classify all finitely generated abelian
groups: as mentioned in UNCERTAIN:Example II.6.3, all such groups are
products of cyclic groups.\footnote{The natural context to prove this
  more general result is that of modules over Euclidean rings or even
  principal ideal domains.} In particular, this is the case for finite
abelian groups: this is what we prove in this section.

Since in this section we exclusively deal with abelian groups, we
revert to the abelian style of notations: thus the operation will be
denoted $+$; the identity will be $0$; direct products will be called
direct sums (and denoted $\oplus$); and so on.

First of all, I will congeal into an explicit statement a simple
observation that has been with us in one form or another since at
least as far back as\footnote{In fact, this observation will really
  find its most appropriate resting place when we prove the Chinese
  Remainder Theorem, \Cref{thm:irred:chinese-remainder-theorem}.}
\Cref{exer:groups:cmn-iso-cm-times-cn}.
\begin{lem}{}{lem:groups2:coprime-order-subgroups-direct-sum}
  Let $G$ be an abelian group, and let $H$, $K$ be subgroups such that
  $|H|$, $|K|$ are relatively prime. Then $H + K \cong H \oplus K$.
\end{lem}
\begin{proof}
  By Lagrange's theorem (\Cref{coro:groups:lagrange}),
  $H \cap K = \{0\}$. Since subgroups of abelian groups are
  automatically normal, the statement follows from
  \cref{prop:groups2:disjoint-normal-subgroups-product-is-direct}.
\end{proof}
Now let $G$ be a \emph{finite} abelian group. For each prime $p$, the
$p$-Sylow subgroup of $G$ is unique, since it is automatically normal
in $G$. Since the distinct nontrivial Sylow subgroups of $G$ are
$p$-groups for different primes $p$, \cref{lem:groups2:coprime-order-subgroups-direct-sum}
immediately implies the following result.

\begin{coro}{}{coro:groups2:finite-abelian-is-sum-of-sylow-subgroups}
  Every finite abelian group is the direct sum of its nontrivial Sylow
  subgroups.
\end{coro}
(The diligent reader knew already that this had to be the case, since
abelian groups are nilpotent; cf.\
\cref{exer:groups2:all-sylow-normal-implies-direct-product}.)  Thus,
we already know that every finite abelian group is a direct sum of
$p$-groups, and our main task amounts to classifying abelian
$p$-groups for a fixed prime $p$. This is somewhat technical; we will
get there by a seemingly roundabout path.
\begin{lem}{}{lem:groups2:max-order-element-splits-sequence}
  Let $G$ be an abelian $p$-group, and let $g \in G$ be an element of
  maximal order. Then the exact sequence
  \[0 \longrightarrow \langle g \rangle \longrightarrow G
    \longrightarrow G/\langle g \rangle \longrightarrow 0\] splits.
\end{lem}

Put otherwise, there is a subgroup $L$ of $G$ such that $L$ maps
isomorphically to $G/\langle g \rangle$ via the canonical projection,
that is, such that $\langle g \rangle \cap L = \{0\}$ and
$\langle g \rangle + L = G$. Note that it will follow that
$G \cong \langle g \rangle \oplus L$, by \cref{prop:groups2:disjoint-normal-subgroups-product-is-direct}.  The
main technicality needed in order to prove this lemma is the following
particular case:

\begin{lem}{}{lem:groups2:noncyclic-abelian-has-extra-order-p-element}
  Let $p$ be a prime integer and $r \geq 1$. Let $G$ be a noncyclic
  abelian group of order $p^{r+1}$, and let $g \in G$ be an element of
  order $p^r$. Then there exists an element $h \in G$,
  $h \notin \langle g \rangle$, such that $|h| = p$.
\end{lem}
\cref{lem:groups2:noncyclic-abelian-has-extra-order-p-element} is a special case of \cref{lem:groups2:max-order-element-splits-sequence} in
the sense that, with notation as in the statement, necessarily
$\langle h \rangle \cong G/\langle g \rangle$, and in fact
$G \cong \langle g \rangle \oplus \langle h \rangle$ (and the reader
is warmly encouraged to understand this before proceeding!). That is,
we can split off the `large' cyclic subgroup $\langle g \rangle$ as a
direct summand of $G$, provided that $G$ is not cyclic and not much
larger than $\langle g \rangle$. \cref{lem:groups2:max-order-element-splits-sequence} claims that
this can be done whenever $\langle g \rangle$ is a maximal cyclic
subgroup of $G$. We will be able to prove this more general statement
easily once the particular case is settled.
\begin{proof}[Proof of \cref{lem:groups2:noncyclic-abelian-has-extra-order-p-element}]
  Denote $\langle g \rangle$ by $K$, and let $h'$ be any element of
  $G$, $h' \notin K$. The subgroup $K$ is normal in $G$ since $G$ is
  abelian; the quotient group $G/K$ has order $p$. Since
  $h' \notin K$, the coset $h' + K$ has order $p$ in $G/K$; that is,
  $ph' \in K$. Let $k = ph'$.

  Note that $|k|$ divides $p^r$; hence it is a power of $p$. Also
  $|k| \neq p^r$, otherwise $|h'| = p^{r+1}$ and $G$ would be cyclic,
  contrary to the hypothesis.

  Therefore $|k| = p^s$ for some $s < r$; $k$ generates a subgroup
  $\langle k \rangle$ of the cyclic group $K$, of order $p^s$. By
  \Cref{prop:groups:subgroups-of-znz},
  $\langle k \rangle = \langle p^{r-s} g \rangle$.  Since $s < r$,
  $\langle k \rangle \subseteq \langle pg \rangle$; thus, $k = mpg$
  for some $m \in \Z$.

  Then let $h = h' - mg$: $h \neq 0$ (since $h' \notin K$), and
  \[ph = ph' - p(mg) = k - k = 0,\] showing that $|h| = p$, as stated.
\end{proof}
\begin{proof}[Proof of \cref{lem:groups2:max-order-element-splits-sequence}]
  Argue by induction on the order of $G$; the case $|G| = p^0 = 1$
  requires no proof. Thus we will assume that $G$ is nontrivial and
  that the statement is true for every $p$-group smaller than $G$.

  Let $g \in G$ be an element of maximal order, say $p^r$, and denote
  by $K$ the subgroup $\langle g \rangle$ generated by $g$; this
  subgroup is normal, as $G$ is abelian. If $G = K$, then the
  statement holds trivially. If not, $G/K$ is a nontrivial $p$-group,
  and hence it contains an element of order $p$ by Cauchy's theorem
  (\cref{thm:groups2:cauchy}). This element generates a subgroup of
  order $p$ in $G/K$, corresponding to a subgroup $G'$ of $G$ of order
  $p^{r+1}$, containing $K$. This subgroup is not cyclic (otherwise
  the order of $g$ is not maximal).

  That is, we are in the situation of \cref{lem:groups2:noncyclic-abelian-has-extra-order-p-element}: hence we
  can conclude that there is an element $h \in G'$ (and hence
  $h \in G$) with $h \notin K$ and $|h| = p$. Let
  $H = \langle h \rangle \subseteq G$ be the subgroup generated by
  $h$, and note that $K \cap H = \{0\}$.

  Now work modulo $H$. The quotient group $G/H$ has smaller size than
  $G$, and $g + H$ generates a cyclic subgroup
  $K' = (K + H)/H \cong K/(K \cap H) \cong K$ of maximal order in
  $G/H$. By the induction hypothesis, there is a subgroup $L'$ of
  $G/H$ such that $K' + L' = G/H$ and $K' \cap L' = \{0_{G/H}\}$. This
  subgroup $L'$ corresponds to a subgroup $L$ of $G$ containing $H$.

  Now I claim that (i) $K + L = G$ and (ii) $K \cap L =
  \{0\}$. Indeed, we have the following:
  \begin{enumerate}[label=(\roman*)]
  \item For any $a \in G$, there exist $mg + H \in K'$,
    $\ell + H \in L'$ such that $a + H = mg + \ell + H$ (since
    $K' + L' = G/H$). This implies $a - mg \in L$, and hence
    $a \in K + L$ as needed.
  \item If $a \in K \cap L$, then
    $a + H \in K' \cap L' = \{0_{G/H}\}$, and hence $a \in H$. In
    particular, $a \in K \cap H = \{0\}$, forcing $a = 0$, as needed.
  \end{enumerate}
  (i) and (ii) imply the lemma, as observed in the comments following
  the statement.
\end{proof}
Now we are ready to state the classification theorem; the proof is
quite straightforward after all this preparation work. We first give
the statement in a somewhat coarse form, as a corollary of the
previous considerations:

\begin{coro}{}{coro:groups2:finite-abelian-is-sum-of-cyclic-p-groups}
  Let $G$ be a finite abelian group. Then $G$ is a direct sum of
  cyclic groups, which may be assumed to be cyclic $p$-groups.
\end{coro}
\begin{proof}
  As noted in \cref{coro:groups2:finite-abelian-is-sum-of-sylow-subgroups}, $G$ is a direct sum of
  $p$-groups (as a consequence of the Sylow theorems). I claim that
  every abelian $p$-group $P$ is a direct sum of cyclic $p$-groups.

  To establish this, argue by induction on $|P|$. There is nothing to
  prove if $P$ is trivial. If $P$ is not trivial, let $g$ be an
  element of $P$ of maximal order. By \cref{lem:groups2:max-order-element-splits-sequence}
  \[P = \langle g \rangle \oplus P'\] for some subgroup $P'$ of $P$;
  by the induction hypothesis $P'$ is a direct sum of cyclic
  $p$-groups, concluding the proof.
\end{proof}

\subsection{Invariant factors and elementary divisors}
\label{subsec:groups2:invariant-factors-and-elementary-divisors}
Here is a more precise version of the classification theorem. It is
common to state the result in two equivalent forms.

\begin{thm}{}{thm:groups2:fundamental-theorem-of-finite-abelian-groups}
  Let $G$ be a finite nontrivial abelian group. Then
  \begin{itemize}
  \item there exist prime integers $p_1, \dots, p_r$ and positive
    integers $n_{ij}$ such that $|G| = \prod_{i,j} p_i^{n_{ij}}$ and
    \[G \cong \bigoplus_{i,j} \frac{\Z}{p_i^{n_{ij}}\Z}\,;\]
  \item there exist positive integers $1 < d_1 \mid \cdots \mid d_s$
    such that $|G| = d_1 \cdots d_s$ and
    \[G \cong \frac{\Z}{d_1\Z} \oplus \cdots \oplus
      \frac{\Z}{d_s\Z}\,.\]
  \end{itemize}
  Further, these decompositions are uniquely determined by $G$.
\end{thm}
The first form is nothing but a more explicit version of the statement
of \cref{coro:groups2:finite-abelian-is-sum-of-cyclic-p-groups}, so it has already been proven. I will
explain how to obtain the second form from the first. The uniqueness
statement\footnote{Of course the `uniqueness' statement only holds up
  to trivial manipulation such as a permutation of the factors. The
  claim is that the factors themselves are determined by $G$, in the
  sense that two direct sums of either form given in the statement are
  isomorphic only if their factors match.} is left to the reader
(\cref{exer:groups2:uniqueness-of-p-grp-decomp}).  The prime powers
appearing in the first form of \cref{thm:groups2:fundamental-theorem-of-finite-abelian-groups} are called the
\emph{elementary divisors} of $G$; the integers $d_i$ appearing in the
second form are called \emph{invariant factors}. To go from elementary
divisors to invariant factors, collect the elementary divisors in a
table, listing (for example) prime powers according to increasing
primes in the horizontal direction and decreasing exponents in the
vertical direction; then the invariant factors are obtained as
products of the factors in each row:
\[
  \begin{array}{c||c|c|c|c}
    d_r     & p_1^{n_{11}} & p_2^{n_{21}} & p_3^{n_{31}} & \cdots \\
    d_{r-1} & p_1^{n_{12}} & p_2^{n_{22}} & p_3^{n_{32}} & \cdots \\
    d_{r-2} & p_1^{n_{13}} & p_2^{n_{23}} & p_3^{n_{33}} & \cdots \\
    \cdots  & \cdots       & \cdots       & \cdots       & \cdots
  \end{array}
\]
Conversely, given the invariant factors $d_i$, obtain the rows of this
table by factoring $d_i$ into prime powers: the condition
$d_1 \mid \cdots \mid d_r$ guarantees that these will be decreasing.

Repeated applications of \cref{lem:groups2:coprime-order-subgroups-direct-sum} show that if
$d = p_1^{n_1} \cdots p_r^{n_r}$ for distinct primes $p_i$ and
positive $n_i$ (as is the case in each row of the table), then
\[\frac{\Z}{d\Z} \cong
  \frac{\Z}{p_1^{n_1}\Z} \oplus \cdots \oplus
  \frac{\Z}{p_r^{n_r}\Z},\] proving that the two decompositions given
in \cref{thm:groups2:fundamental-theorem-of-finite-abelian-groups} are indeed equivalent. This will likely be
much clearer once the reader works through a few examples.
\begin{exam}{}{exam:groups2:decomposition-of-group-of-order-29160}
  Here are the two decompositions for a (random) group of order
  $29160 = 2^3 \cdot 3^6 \cdot 5$:
  \begin{align*}
    &\frac{\Z}{2\Z} \oplus
      \frac{\Z}{2\Z} \oplus
      \frac{\Z}{2\Z} \oplus
      \frac{\Z}{3\Z} \oplus
      \frac{\Z}{3\Z} \oplus
      \frac{\Z}{3^2\Z} \oplus
      \frac{\Z}{3^2\Z} \oplus
      \frac{\Z}{5\Z} \\
    \cong\,&\frac{\Z}{3\Z} \oplus
             \frac{\Z}{6\Z} \oplus
             \frac{\Z}{18\Z} \oplus
             \frac{\Z}{90\Z}
  \end{align*}
  and here is the corresponding table of invariant factors/elementary
  divisors:
  \[
    \begin{array}{c||c|c|c}
      90 & 2 & 3^2 & 5 \\
      18 & 2 & 3^2 &   \\
      6 & 2 & 3   &   \\
      3 &   & 3   &
    \end{array}
  \]
\end{exam}
\begin{exam}{}{exam:groups2:abelian-groups-of-order-360}
  There are exactly $6$ isomorphism classes of abelian groups of order
  $360$. Indeed, $360 = 2^3 \cdot 3^2 \cdot 5$; the six possible
  tables of elementary divisors are shown below. In terms of invariant
  factors, the six distinct abelian groups of order $360$ (up to
  isomorphism, by the uniqueness part of \cref{thm:groups2:fundamental-theorem-of-finite-abelian-groups}) are
  therefore
  \begin{gather*}
    \frac{\Z}{360\Z},\quad
    \frac{\Z}{2\Z} \oplus
    \frac{\Z}{180\Z},\quad
    \frac{\Z}{2\Z} \oplus
    \frac{\Z}{2\Z} \oplus
    \frac{\Z}{90\Z},\\
    \frac{\Z}{3\Z} \oplus
    \frac{\Z}{120\Z},\quad
    \frac{\Z}{6\Z} \oplus
    \frac{\Z}{60\Z},\quad
    \frac{\Z}{2\Z} \oplus
    \frac{\Z}{6\Z} \oplus
    \frac{\Z}{30\Z}.
  \end{gather*}
  \[
    \begin{array}{c||c|c|c}
      360 & 2^3 & 3^2 & 5
    \end{array}
    \quad
    \begin{array}{c||c|c|c}
      180 & 2^2 & 3^2 & 5 \\
      2 & 2   &     &
    \end{array}
    \quad
    \begin{array}{c||c|c|c}
      90 & 2 & 3^2 & 5 \\
      2 & 2 &     &   \\
      2 & 2 &     &
    \end{array}
  \]
  \[
    \begin{array}{c||c|c|c}
      120 & 2^3 & 3 & 5 \\
      3 &     & 3 &
    \end{array}
    \quad
    \begin{array}{c||c|c|c}
      60 & 2^2 & 3 & 5 \\
      6 & 2   & 3 &
    \end{array}
    \quad
    \begin{array}{c||c|c|c}
      30 & 2 & 3 & 5 \\
      6 & 2 & 3 &   \\
      2 & 2 &   &
    \end{array}
  \]
\end{exam}

\subsection{Application: Finite subgroups of multiplicative groups of
  fields}
\label{subsec:groups2:application-finite-subgroups-of-multiplicative-groups-of-fields}
Any classification theorem is useful in that it potentially reduces
the proof of general facts to explicit verifications. Here is one
example illustrating this strategy:
\begin{lem}{}{lem:groups2:finite-group-with-bounded-n-torsion-is-cyclic}
  Let $G$ be a finite abelian group, and assume that for every integer
  $n > 0$ the number of elements $g \in G$ such that $ng = 0$ is at
  most $n$. Then $G$ is cyclic.
\end{lem}
The reader should try to prove this `by hand', to appreciate the fact
that it is not entirely trivial. It does become essentially immediate
once we take the classification of finite abelian groups into account.
Indeed, by \cref{thm:groups2:fundamental-theorem-of-finite-abelian-groups}
\[G \cong \frac{\Z}{d_1\Z} \oplus \cdots \oplus \frac{\Z}{d_s\Z}\] for
some positive integers $1 < d_1 \mid \cdots \mid d_s$. But if $s > 1$,
then $|G| > d_s$ and $d_s g = 0$ for all $g \in G$ (so that the order
of $g$ divides $d_s$), contradicting the hypothesis.  Therefore
$s = 1$; that is, $G$ is cyclic.  \cref{lem:groups2:finite-group-with-bounded-n-torsion-is-cyclic} is the key to
a particularly nice proof of the following important fact, a weak form
of which\footnote{The diligent reader has proved that particular case
  in \Cref{exer:groups:zpz-star-cyclic}. The proof hinted at in that
  exercise upgrades easily to the general case presented here. The
  point is not that the classification theorem is \emph{necessary} in
  order to prove statements such as \cref{thm:groups2:finite-subgroup-of-field-mult-group-is-cyclic}; the point
  is that it makes such statements nearly evident.} we ran across back
in \Cref{exam:groups:cyclic-formal-intro}. Recall that the set $F^*$
of nonzero elements of a field $F$ is a commutative group under
multiplication. Also recall (\Cref{exam:rings:remainder-degree-1})
that a polynomial $f(x) \in F[x]$ is divisible by $(x - a)$ if and
only if $f(a) = 0$; since a nonzero polynomial of degree $n$ over a
field can have at most $n$ linear factors, this shows\footnote{Unique
  factorization in $F[x]$ is secretly needed here. We will deal with
  this issue more formally later; cf.\ \Cref{lem:irred:polynomial-root-count-bound}.} that if
$f(x) \in F[x]$ has degree $n$, then $f(a) = 0$ for at most $n$
distinct elements $a \in F$.
\begin{thm}{}{thm:groups2:finite-subgroup-of-field-mult-group-is-cyclic}
  Let $F$ be a field, and let $G$ be a finite subgroup of the
  multiplicative group $(F^*, \cdot)$. Then $G$ is cyclic.
\end{thm}
\begin{proof}
  By the considerations preceding the statement, for every $n$ there
  are at most $n$ elements $a \in F$ such that $a^n - 1 = 0$, that is,
  at most $n$ elements $a \in G$ such that $a^n = 1$.
  \cref{lem:groups2:finite-group-with-bounded-n-torsion-is-cyclic} implies then that $G$ is cyclic.
\end{proof}

As a (very) particular case, the multiplicative group
$((\Z/p\Z)^*, \cdot)$ is cyclic: this is the fact pointed to in
\Cref{exam:groups:cyclic-formal-intro}.  \emph{Preview of coming
  attractions:} Finitely generated (as opposed to just finite) abelian
groups are also direct sums of cyclic groups.  The only difference
between the classification of finitely generated abelian groups and
the classification of finite abelian groups explored here is the
possible presence of a `free' factor $\Z^{\oplus r}$ in the
decomposition. The reader will first prove this fact in VI.2.19, as a
consequence of `Gaussian elimination over integral domains', and then
recover it again as a particular case of the classification theorem
for finitely generated modules over PIDs, \Cref{thm:linalg:structure-theorem-fg-modules-over-pid}. Neither
Gaussian elimination nor the very general \Cref{thm:linalg:structure-theorem-fg-modules-over-pid} are any
harder to prove than the particular case of finite abelian groups
laboriously worked out by hand in this section---a common benefit of
finding the right general point of view is that, as a rule, proofs
simplify. Technical work such as that performed in order to prove
\cref{lem:groups2:max-order-element-splits-sequence} is absorbed into the work necessary to build up
the more general apparatus; the need for such technicalities
evaporates in the process.

\subsection*{Exercises}

\begin{exer}{}{exer:groups2:uniqueness-of-p-grp-decomp}
  \exerused{} Prove that the decomposition of a finite abelian group
  $G$ as a direct sum of cyclic $p$-groups is unique. (Hint: The prime
  factorization of $|G|$ determines the primes, so it suffices to show
  that if
  \[\frac{\Z}{p^{r_1}\Z} \oplus \dots \oplus
    \frac{\Z}{p^{r_m}\Z} \cong \frac{\Z}{p^{s_1}\Z} \oplus \dots
    \oplus \frac{\Z}{p^{s_n}\Z},\] with $r_1 \geq \dots \geq r_m$ and
  $s_1 \geq \dots \geq s_n$, then $m = n$ and $r_i = s_i$ for all
  $i$. Do this by induction, by considering the group $pG$ obtained as
  the image of the homomorphism $G \to G$ defined by $g \mapsto pg$.)
  [\Cref{subsec:groups2:invariant-factors-and-elementary-divisors},
  \Cref{subsec:linalg:the-classification-theorem},
  \ref{exer:linalg:complete-uniqueness-proof-classification-thm}]
\end{exer}
\begin{exer}{}{exer:groups2:complete-classification-ord-8}
  Complete the classification of groups of order $8$ (cf.\
  \cref{exer:groups2:noncomm-grp-ord-8-is-d8-or-q8}).
\end{exer}
\begin{exer}{}{exer:groups2:noncomm-grp-ord-p3-ctr-and-quot}
  Let $G$ be a noncommutative group of order $p^3$, where $p$ is a
  prime integer. Prove that $Z(G) \cong \Z/p\Z$ and
  $G/Z(G) \cong \Z/p\Z \times \Z/p\Z$.
\end{exer}
\begin{exer}{}{exer:groups2:classify-abel-grps-ord-400}
  Classify abelian groups of order $400$.
\end{exer}
\begin{exer}{}{exer:groups2:num-abel-grps-ord-pr-eq-partitions}
  Let $p$ be a prime integer. Prove that the number of distinct
  isomorphism classes of abelian groups of order $p^r$ equals the
  number of partitions of the integer $r$.
\end{exer}
\begin{exer}{}{exer:groups2:num-abel-grps-ord-1024}
  \exerused{} How many \emph{abelian} groups of order $1024$ are
  there, up to isomorphism?
  [\Cref{subsec:groups:example-subgroup-generated-by-a-subset}]
\end{exer}
\begin{exer}{}{exer:groups2:mult-by-p-hom-ker-coker-props}
  $\lnot$ Let $p > 0$ be a prime integer, $G$ a finite abelian group,
  and denote by $\rho : G \to G$ the homomorphism defined by
  $\rho(g) = pg$.
  \begin{itemize}
  \item Let $A$ be a finite abelian group such that $pA = 0$. Prove
    that $A \cong \Z/p\Z \oplus \dots \oplus \Z/p\Z$.
  \item Prove that $p\ker\rho$ and $p(\coker\rho)$ are both $0$.
  \item Prove that $\ker\rho \cong \coker\rho$.
  \item Prove that every subgroup of $G$ of order $p$ is contained in
    $\ker\rho$ and that every subgroup of $G$ of index $p$ contains
    $\image\rho$.
  \item Prove that the number of subgroups of $G$ of order $p$ equals
    the number of subgroups of $G$ of index $p$.
  \end{itemize}
  [\cref{exer:groups2:subgrp-with-given-invariant-divisors}]
\end{exer}
\begin{exer}{}{exer:groups2:subgrp-with-given-invariant-divisors}
  $\lnot$ Let $G$ be a finite abelian $p$-group, with elementary
  divisors $p^{n_1}, \dots, p^{n_r}$ ($n_1 \geq n_2 \geq \dots$).
  Prove that $G$ has a subgroup $H$ with elementary divisors
  $p^{m_1}, \dots, p^{m_s}$ ($m_1 \geq m_2 \geq \dots$) if and only if
  $s \leq r$ and $m_i \leq n_i$ for $i = 1, \dots, s$. (Hint: One
  direction is immediate. For the other, with notation as in
  \cref{exer:groups2:mult-by-p-hom-ker-coker-props}, compare
  $\ker\rho$ for $H$ and $G$ to establish $s \leq r$; this also proves
  the statement if all $n_i = 1$.  For the general case use induction,
  noting that if $G \cong \bigoplus_i \Z/p^{n_i}\Z$, then
  $\rho(G) \cong \bigoplus_i \Z/p^{n_i - 1}\Z$.)

  Prove that the same description holds for the homomorphic images of
  $G$. [\cref{exer:groups2:fin-abel-grp-contains-subgrp-iso-quot}]
\end{exer}
\begin{exer}{}{exer:groups2:fin-abel-grp-contains-subgrp-iso-quot}
  Let $H$ be a subgroup of a finite abelian group $G$. Prove that $G$
  contains a subgroup isomorphic to $G/H$. (Reduce to the case of
  $p$-groups; then use
  \cref{exer:groups2:subgrp-with-given-invariant-divisors}.) Show that
  both hypotheses `finite' and `abelian' are needed for this
  result. (Hint: $Q_8$ has a unique subgroup of order $2$.)
\end{exer}
\begin{exer}{}{exer:groups2:dual-of-fin-abel-grp-iso-self}
  The \emph{dual} of a finite group $G$ is the abelian group
  $G^\vee := \Hom_{\textsf{Grp}}(G, \C^*)$, where $\C^*$ is the
  multiplicative group of $\C$.
  \begin{itemize}
  \item Prove that the image of every $\sigma \in G^\vee$ consists of
    \emph{roots of $1$} in $\C$, that is, roots of polynomials
    $x^n - 1$ for some $n$.
  \item Prove that if $G$ is a finite abelian group, then
    $G \cong G^\vee$. (Hint: First prove this for cyclic groups; then
    use the classification theorem to generalize to the arbitrary
    case.)
  \end{itemize}
  In~\Cref{subsec:linrep:ext-Z-G-Z} we will encounter another
  notion of `dual' of a group.
\end{exer}
\begin{exer}{}{exer:groups2:classify-modules-over-znz}
  \begin{itemize}
  \item Use the classification theorem for finite abelian groups
    (\cref{thm:groups2:fundamental-theorem-of-finite-abelian-groups}) to classify all finite modules over the
    ring $\Z/n\Z$.
  \item Prove that if $p$ is prime, all finite modules over $\Z/p\Z$
    are free.\footnote{As we will see in \Cref{prop:linalg:field-iff-fg-modules-free}, this
      property characterizes fields.}
  \end{itemize}
\end{exer}
\begin{exer}{}{exer:groups2:cancellation-for-fin-abel-grps}
  Let $G$, $H$, $K$ be finite abelian groups such that
  $G \oplus H \cong G \oplus K$. Prove that $H \cong K$.
\end{exer}
\begin{exer}{}{exer:groups2:same-elt-counts-implies-iso-abel}
  $\lnot$ Let $G$, $H$ be finite abelian groups such that, for all
  positive integers $n$, $G$ and $H$ have the same number of elements
  of order $n$. Prove that $G \cong H$. (Note: The `abelian'
  hypothesis is necessary! $C_4 \times C_4$ and $Q_8 \times C_2$ are
  nonisomorphic groups both with $1$ element of order $1$, $3$
  elements of order $2$, and $12$ elements of order $4$.)
  [\Cref{subsec:groups:isomorphisms}]
\end{exer}
\begin{exer}{}{exer:groups2:unique-subgrp-ord-p-implies-cyclic}
  Let $G$ be a finite abelian $p$-group, and assume $G$ has only one
  subgroup of order $p$. Prove that $G$ is cyclic. (This is in some
  sense a converse to \Cref{prop:groups:subgroups-of-znz}. You are
  welcome to try to prove it `by hand', but use of the classification
  theorem will simplify the argument considerably.)
\end{exer}
\begin{exer}{}{exer:groups2:ord-of-elts-divides-max-ord}
  Let $G$ be a finite abelian group, and let $a \in G$ be an element
  of \emph{maximal} order in $G$. Prove that the order of every
  $b \in G$ divides $|a|$. (This essentially reproduces the result of
  \Cref{exer:groups:maximal-order-divides}.)
\end{exer}
\begin{exer}{}{exer:groups2:at-most-one-subgrp-per-divisor-implies-cyclic}
  Let $G$ be an abelian group of order $n$, and assume that $G$ has at
  most one subgroup of order $d$ for all $d \mid n$. Prove that $G$ is
  cyclic.
\end{exer}

%%% Local Variables:
%%% mode: LaTeX
%%% TeX-engine: luatex
%%% TeX-master: "../master"
%%% End:
